% !TEX encoding = UTF-8 Unicode
%% NOTE: Due to using fontspec and UTF8 to avoid encoding problems and stuff, Russian
%% version of this document builds using xe(la)tex and TexStudio
\documentclass[russian,14pt,oneside]{book}
\usepackage{extsizes} % additional font sizes for documentclass

%%%%%%
%% package includes
\usepackage{ifxetex}
\ifxetex
\else
	\usepackage[utf8]{inputenc}
	\usepackage[T2A]{fontenc}
\fi

\usepackage[russian]{babel}
\usepackage{hyperref}
\usepackage{textcomp}
\usepackage{array}
\usepackage{color}
\usepackage[hmargin=2.5cm,vmargin=2.5cm]{geometry}
\usepackage{minted}
\usepackage{todonotes}
\usepackage{parskip}
\usepackage[Bjornstrup]{fncychap}
\usepackage{listings}

%%%%%%
%% package setup
\hypersetup{
    colorlinks=true,
    linktoc=all,
    linkcolor=red,
}

%%%%%%
%% Fonts
\ifxetex
    \usepackage{fontspec}
    \defaultfontfeatures{Scale=MatchLowercase,Mapping=tex-text}
    \setmainfont[Path=fonts/,ItalicFont=PTF56F.ttf,
        BoldFont=PTF75F.ttf]{PTF55F.ttf}
    \setmonofont[Ligatures={NoRequired,NoCommon,NoContextual},
        Path=fonts/,BoldFont=FiraMono-Bold.ttf]{FiraMono-Regular.ttf}
    \usepackage{titlesec}
    \newfontfamily\erlChapterFont[Path=fonts/,BoldFont=PTF75F.ttf]{PTF55F.ttf}
    \ChTitleVar{\erlChapterFont\fontsize{28}{28}\bf\selectfont\hfill}
    \ChNameVar{\erlChapterFont\fontsize{16}{16}\bf\selectfont}
\fi

%%%%%%
%% hacks and macros
\newminted[erlang]{erlang}{xleftmargin=0cm,xrightmargin=0cm,frame=single,framerule=1pt}

%% Framed code fragment coloring, for code examples where Russian variable 
%% names break syntax and coloring
\usepackage{fancyvrb}
\DefineVerbatimEnvironment{erlangru}{Verbatim}{frame=single,samepage=true}

%% hack to fix the margin issue in \item s in 1.2.1
\newminted[erlangim]{erlang}{xleftmargin=-1.5cm}

%%% display todo notes or not?
\newcommand{\todonote}[1]{\todo[inline]{\textbf{Note:} #1}}
%\newcommand{\todonote}[1]{}

%% typesetting for the "results in" notation used.
\def\resultingin{$\quad\Rightarrow\quad$}

\usepackage{wrapfig}

\begin{document}

%%%%%%
%% Title
\begin{titlepage}
\centering

\vspace*{70pt}
\includegraphics[scale=0.3]{includes/erlang-logo.png}\\[0.8\baselineskip]
{\Huge \sffamily НАСТОЛЬНАЯ КНИГА}\\
\vspace{250pt}
{\LARGE Бьярн Дэкер (Bjarne D\"acker)}\par
{\LARGE Роберт Вирдинг (Robert Virding)}\par

\end{titlepage}


%%%%%%
%% Info Page
\clearpage
\thispagestyle{empty}
{\Huge Настольная книга по Erlang}\\[0.1\baselineskip]
\hspace{10pt}{\large{Бьярн Дэкер (Bjarne Däcker) и Роберт Вирдинг (Robert Virding)}}

\vspace{10pt}
{\Large \textbf{Версия:}\\[0.2\baselineskip]
\immediate\write18{git log -n1 | grep 'Date:' | sed 's/Date:   //g' > _revision.tex}
\input{_revision}
\immediate\write18{rm _revision.tex}
}

\vspace{10pt}
{\large Последняя версия этой книги находится по адресу:\\
\url{http://opensource.erlang-solutions.com/erlang-handbook}}

\vfill

\textbf{Главный редактор}\\[0.1\baselineskip]
Омер Килич (Omer Kilic)

\textbf{Прочие участники}\\[0.1\baselineskip]
Перевод на русский: Дмитрий Литовченко (Dmytro Lytovchenko)

Список людей, которые помогли с изменениями и исправлениями находится
\href{https://github.com/esl/erlang-handbook/graphs/contributors}{в репозитории 
проекта}.

\vspace{10pt}
\textbf{Соглашения}\\
Спецификации синтаксиса задаются с помощью \texttt{этого моноширинного шрифта}.
В квадратные скобки ([ ]) заключаются необязательные части. Термы, начинающиеся
с заглавной буквы, как например \textit{Integer} должны быть заменены на 
какое-нибудь подходящее значение. Термы, начинающиеся со строчной буквы, как
например \texttt{end}, являются зарезервированными словами языка Erlang. 
Вертикальная черта (\textbar{}) разделяет альтернативные варианты, как например 
Integer \textbar{} Float.

\vspace{10pt}
\textbf{Ошибки и улучшения}\\
Это живой документ, пожалуйста присылайте исправления и рекомендации по улучшению
содержимого используя систему учёта задач Github по адресу
\url{https://github.com/esl/erlang-handbook}. Вы также можете создать свою ветвь
репозитория (fork) и прислать нам запрос на соединение ветвей (pull request) с
вашими предлагаемыми исправлениями и предложениями. Новые ревизии этого документа
будут публиковаться по мере накопления больших изменений.

\vspace{10pt}
\begin{wrapfigure}[2]{l}{0.15\textwidth}
	\vspace{-12pt}
	\includegraphics[scale=0.7]{includes/cc-by-sa.png}
	\vspace{-12pt}
\end{wrapfigure}
Этот текст доступен согласно условиям лицензии Creative Commons
Attribution-ShareAlike 3.0 License. Вы имеете право копировать, распространять и
передавать эту книгу по условиям лицензии,  описанным по адресу
\url{http://creativecommons.org/licenses/by-sa/3.0}

\newpage


%%%%%%
%% TOC
\tableofcontents


%%%%%%
%% Chapters

\chapter[Вступление, или почему Erlang такой, как он есть]{Вступление}
\section*{или почему Erlang такой, как он есть}
\label{background}

Erlang --- это результат проекта по улучшению программирования приложений для 
телекоммуникаций в Лаборатории Компьютерных Наук (CSLab или Computer Science Lab)
компании Ericsson. Критически важным требованием была поддержка характеристик 
этих приложений, таких, как: 

\begin{itemize}
	\item Массивная параллельность
	\item Устойчивость к сбоям
	\item Изоляция
	\item Динамическое обновление кода во время его исполнения
	\item Транзакционность
\end{itemize}

В течение всей истории Erlang процесс его разработки был исключительно 
прагматичным. Характеристики и свойства видов систем, в которых была заинтересована
компания Ericsson, прямым образом влияли на ход разработки Erlang. Эти свойства
считались настолько фундаментальными, что было решено поддержку для них встроить
прямо в язык, вместо дополнительных библиотек. По причине прагматичного процесса
разработки вместо предварительного планирования, Erlang "<стал"> функциональным
языком --- поскольку свойства функциональных языков очень хорошо подходили к
свойствам систем, которые разрабатывались.


\chapter{Структура Erlang-программы}

\section{Синтаксис модулей}

Программа на Erlang состоит из \textbf{модулей} где каждый модуль является текстовым
файлом с расширением \textbf{.erl}. Для небольших программ все модули обычно хранятся
в одной директории. Модуль состоит из атрибутов модуля и определений функций.

\begin{erlang}
-module(demo).
-export([double/1]).

double(X) -> times(X, 2).

times(X, N) -> X * N.
\end{erlang}

Показанный модуль \texttt{demo} состоит из функции \texttt{times/2}, которая является
локальной для модуля и функции \texttt{double/1}, которая экспортирована и может быть
вызвана снаружи модуля.

\texttt{demo:double(10)} \resultingin \texttt{20}\hfill
(стрелка $\Rightarrow$ читается как "<даёт результат">)

\texttt{double/1} означает, что перед нами функция "<double"> с \textit{одним}
аргументом. Функция \texttt{double/2} принимает \textit{два} аргумента и считается
другой,  отличной от первой, функцией. Количество аргументов называется
\textbf{арностью} (arity) данной функции.



\section{Атрибуты модулей}
\textbf{Атрибут модуля} определяет некоторое свойство модуля и состоит из 
\textbf{метки} и \textbf{значения} в следующем виде:

\texttt{-Метка(Значение).}

\texttt{Метка} должна быть атомом, в то время, как значение \texttt{Значение} 
должно быть термом (смотрите главу \ref{datatypes}). Можно указывать любое имя
для атрибута модуля, также разрешены повторы. Атрибуты хранятся в скомпилированном
коде и могут быть извлечены с помощью вызова функции
\texttt{Модуль:module\_info(attributes).}

\subsection{Предопределённые атрибуты модулей}
Предопределённые атрибуты модулей должны быть размещены в модуле до начала первого
определения функции.

\begin{itemize}

	\item \verb|-module(Модуль).|\hfill\\
	Этот атрибут обязателен и должен идти первой строкой файла. Он определяет имя
	модуля. Имя Module, атом (смотрите раздел \ref{datatypes:atom}),
    должно соответствовать имени файла без расширения \texttt{.erl}.

	\item \verb|-export([Функ1/Арность1, ..., ФункN/АрностьN]).|\hfill\\
	Этот атрибут указывает, какие функции модуля могут быть вызваны снаружи 
	модуля. Каждое имя функции \texttt{Функ\textsubscript{X}} --- это атом и 
	\texttt{АрностьX} функции --- целое число.

	\item \verb|-import(Модуль,[Функ1/Арность1, ..., ФункN/АрностьN]).|\hfill\\
	Этот атрибут указывает модуль \texttt{Модуль}, из которого импортируется список
	функций. Например:

	\begin{itemize}
	\item \verb|import(demo, [double/1]).|\hfill\\
	Эта запись означает, что теперь можно писать \texttt{double(10)} вместо более
	длинной записи \texttt{demo:double(10)}, которая может быть непрактична, если
	функция используется часто.
	\end{itemize}

	\item \verb|-compile(Параметры).|\hfill\\
	Параметры компилятора для данного модуля.

	\item \verb|-vsn(Vsn).|\hfill\\
	Версия модуля. Если не указать этот атрибут, то по умолчанию будет использована
	контрольная сумма содержимого модуля.

	\item \verb|-behaviour(Поведение).|\hfill\\
	Этот атрибут указывает поведение модуля, либо выбранное пользователем, либо 
	одно из стандартных поведений OTP: \texttt{gen\_server}, \texttt{gen\_fsm},
	\texttt{gen\_event} или \texttt{supervisor}. Принимаются как британская 
	(behaviour), так и американская запись (behavior).
\end{itemize}



\subsection{Определения записей и макросов}

Записи и макросы определяются так же, как и другие атрибуты модуля:

\begin{verbatim}
-record(Запись, Поля).
-define(Макро, Замена).
\end{verbatim}

Записи и определения макросов также позволяются между функциями, если определение
встречается раньше, чем его первое использование. (Подробнее о записях читайте 
секцию \ref{datatypes:record}, а о макросах главу \ref{macros}.)



\subsection{Включение содержимого файлов}

Включение содержимого файлов указывается аналогично другим атрибутам модуля:

\begin{verbatim}
-include(Файл).
-include_lib(Файл).
\end{verbatim}

\texttt{File} --- это строка, представляющая собой имя файла. Включаемые файлы 
обычно используются для определений записей и макросов, которые используются
в нескольких модулях сразу. По договорённости для включаемых файлов используется
расширение \texttt{.hrl}.

\begin{erlang}
-include("my_records.hrl").
-include("incdir/my_records.hrl").
-include("/home/user/proj/my_records.hrl").
\end{erlang}

Если путь к файлу \texttt{File} начинается с компонента пути \texttt{\$Var}, то
значение переменной окружения \texttt{Var} (которое возвращается функцией
\texttt{os:getenv(Var)}) будет подставлено вместо \texttt{\$Var}.

\begin{erlang}
-include("$PROJ_ROOT/my_records.hrl").
\end{erlang}
%%$ texmaker parser bug

Директива \texttt{include\_lib} подобна \texttt{include}, но считается, что 
первый компонент пути --- имя приложения..

\begin{erlang}
-include_lib("kernel/include/file.hrl").
\end{erlang}

Сервер кода (специальный процесс, управляющий загрузкой и версиями модулей в памяти) 
использует функцию \texttt{code:lib\_dir(kernel)} для поиска директории текущей
(свежайшей) версии модуля \texttt{kernel}, и затем в поддиректори \texttt{include} 
производится поиск нужного файла \texttt{file.hrl}.



\section{Комментарии}

Комментарии могут появляться в любом месте модуля, кроме как внутри строк и атомов,
заключённых в кавычки. Комментарий начинается с символа процента (\texttt{\%}) и
действителен до конца строки но не включая символ конца строки. Завершающий строку
символ конца строки, с точки зрения компилятора, действует как пробел.



\section{Кодировка файлов}

Erlang работает с полным набором символов Latin-1 (ISO-8859-1). Таким образом, 
можно использовать и выводить на экран все печатные символы из набора Latin-1 без
цитирования с помощью обратной косой черты (\textbackslash{}). Атомы и переменные 
могут использовать все символы из набора Latin-1.

\vspace*{12pt}
\begin{center}
\begin{tabular}{|>{\raggedright}p{52pt}|>{\raggedright}p{53pt}|>{\raggedright}p{103pt}|>{\raggedright}p{87pt}|}
\hline
\multicolumn{4}{|p{297pt}|}{Классы символов в кодировке}\tabularnewline
\hline
Восьме\-ричное & Деся\-тичное~ &   & Класс \tabularnewline
\hline
40 -  57 & 32 - 47 &  ! \texttt{"} \# \$ \% \& ' / & Символы пунктуации \tabularnewline
\hline
60 -  71 & 48 - 57 & 0 - 9 & Десятичные цифры \tabularnewline
\hline
72 - 100 & 58 - 64 & : ; \texttt{<} = \texttt{>} @ & Символы пунктуации \tabularnewline
\hline
101 - 132 &  65 - 90 & A - Z & Заглавные буквы \tabularnewline
\hline
133 - 140 &  91 - 96 & [ \textbackslash{} ] \textasciicircum{} \_ ` & Символы пунктуации \tabularnewline
\hline
141 - 172 &  97 - 122 & a  -  z & Строчные буквы \tabularnewline
\hline
173 - 176 & 123 - 126 & \{ \textbar{} \} \textasciitilde{} & Символы пунктуации \tabularnewline
\hline
200 - 237 & 128 - 159 ~ &   & Управляющие символы \tabularnewline
\hline
240 - 277 & 160 - 191 & - ¿  & Символы пунктуации \tabularnewline
\hline
300 - 326 & 192 - 214 & À - Ö  & Заглавные буквы \tabularnewline
\hline
327  & 215 & ×  & Символ пунктуации \tabularnewline
\hline
330 - 336 & 216 - 222 & Ø - Þ  & Заглавные буквы \tabularnewline
\hline
337 - 366 & 223 - 246 & ß - ö  & Строчные буквы \tabularnewline
\hline
367  & 247 & ÷  & Символ пунктуации \tabularnewline
\hline
370 - 377 & 248 - 255 & ø - ÿ  & Строчные буквы \tabularnewline
\hline
\end{tabular}
\end{center}

% because of where this lands, force a page break to avoid orphan.
\newpage
\section{Зарезервированные слова}

%\vspace{12pt}

Следующие ключевые слова в Erlang зарезервированы, и не могут использоваться в
качестве атомов (для использования в качестве атома должны быть заключены в 
одиночные кавычки):

\begin{erlang}
after and andalso band begin bnot bor bsl bsr bxor case catch cond
div end fun if let not of or orelse receive rem try when xor
\end{erlang}

\chapter{Типы данных (термы)}
\label{datatypes}


\section{Унарные (одиночные) типы данных}


\subsection{Атомы}
\label{datatypes:atom}
\textbf{Атом} --- это символьное имя, также известное, как \emph{литерал}.
Атомы начинаются со строчной латинской буквы и могут содержать буквенно-цифровые
символы, подчёркивания (\texttt{\_}) и символ \emph{at} (\texttt{@}).
Как вариант, атомы могут быть указаны с помощью заключения в одиночные кавычки
(\texttt{'}), это необходимо, если атом начинается с заглавной буквы, или 
содержит другие символы, кроме подчёркиваний и \emph{at} (\texttt{@}). 
Например:

\begin{erlang}
hello
phone_number
'Monday'
'phone number'
\end{erlang}

\texttt{'Anything inside quotes \textbackslash n\textbackslash 012'}
\hfill(см. раздел \ref{datatypes:escapeseq})


\subsection{Истина и ложь}
\label{datatypes:boolean}

В Erlang нет специального типа данных для \textbf{логических} значений. Эту роль
выполняют атомы \texttt{true} и \texttt{false}.

\texttt{2 =< 3} \resultingin \texttt{true} \\
\texttt{true or false} \resultingin \texttt{true}


\subsection{Целые числа}
\label{datatypes:integer}

В дополнение к обычному способу записи \textbf{целых чисел} Erlang предлагает
ещё ряд способов записи. Запись \texttt{\$Char} даёт числовое значение для 
символа '\texttt{Char}' в кодировке Latin-1 (это также может быть и непечатный
символ) и запись\linebreak
\texttt{Base\#Value} --- даёт целое число, записанное с основанием
\texttt{Base}, основание должно быть целым числом в диапазоне $2..36$.

\texttt{42} \resultingin \texttt{42} \\
\$A  \resultingin \texttt{65} \\
\texttt{\$\textbackslash n} \resultingin \texttt{10}\hfill(see section \ref{datatypes:escapeseq}) \\
\texttt{2\#101} \resultingin \texttt{5} \\
\texttt{16\#1f} \resultingin \texttt{31}


\subsection{Действительные с плавающей точкой}
\label{datatypes:float}

Действительное число \textbf{с плавающей точкой} записывается, как\linebreak
\texttt{Основание}[\texttt{"e"Экспонента}], где \texttt{Основание} --- это
действительное число в диапазоне от 0.01 до 10000 и \texttt{Экспонента} 
(необязательное) --- это целое число со знаком, указывающее экспоненту (степень
десятки, на которую множится \texttt{Основание}). Например:

\texttt{2.3e-3} \resultingin \texttt{2.30000e-3}\hfill
(соответствует $2.3 \cdot 10^{-3}$)


\subsection{Ссылочные значения}
\label{datatypes:reference}

\textbf{Ссылка} --- это терм, значение которого является уникальным в системе 
времени исполнения Erlang, который создаётся встроенной функцией 
\texttt{make\_ref/0}. (Для дополнительной информации по встроенным функциям,
или \emph{BIF}-ам, смотрите раздел \ref{functions:bifs}.)


\subsection{Порты}
\label{datatypes:port}

\textbf{Идентификатор порта} указывает на открытый в системе порт (смотрите 
главу \ref{ports} про порты).


\subsection{Идентификаторы процессов (Pid)}
\label{datatypes:pid}

\textbf{Идентификатор процесса}, или \emph{pid}, указывает на процесс в системе 
(смотрите главу \ref{processes} о процессах).


\subsection{Анонимные функции}
\label{datatypes:fun}

Тип данных \emph{fun} идентифицирует \emph{анонимную функцию} или 
\textbf{функциональный объект} (см. раздел \ref{functions:funs}).



\section{Составные типы данных}


\subsection{Кортежи}
\label{datatypes:tuple}

\textbf{Кортеж} --- это составной тип данных, который содержит 
\textbf{фиксированное количество термов}, заключённое \{в фигурные скобки\}.
Например:

\begin{erlangru}
{Терм1,...,ТермN}
\end{erlangru}

Каждый из \texttt{ТермX} в кортеже называется \textbf{элементом}. Количество 
элементов называется \textbf{размером} данного кортежа.

\begin{center}
\begin{tabular}{|>{\raggedright}p{220pt}|>{\raggedright}p{230pt}|}
\hline
\multicolumn{2}{|p{321pt}|}{BIF-функции для работы с кортежами}\tabularnewline
\hline
\texttt{size(Кортеж)} & 
Возвращает размер \texttt{Кортежа}\tabularnewline
\hline
\texttt{element(N,Кортеж)} &
Возвращает \texttt{N}\textsuperscript{ый} элемент в \texttt{Кортеже} 
\tabularnewline
\hline
\texttt{setelement(N,Кортеж1,Выражение)} &
Возвращает новый кортеж, скопированный из \texttt{Кортеж1}, в котором 
\texttt{N}\textsuperscript{ый} элемент заменён новым значением 
\texttt{Выражение}\tabularnewline
\hline
\end{tabular}
\end{center}

\texttt{P = \{adam, 24, \{july, 29\}\}} \resultingin \texttt{P} связано с 
	\texttt{\{adam, 24, \{july, 29\}\}} \\
\texttt{element(1, P)} \resultingin \texttt{adam} \\
\texttt{element(3, P)} \resultingin \texttt{ \{july,29\}} \\
\texttt{P2 = setelement(2, P, 25)} \resultingin \texttt{P2} связано с 
	\texttt{\{adam, 25, \{july, 29\}\}} \\
\texttt{size(P)} \resultingin \texttt{3} \\
\texttt{size(\{\})} \resultingin \texttt{0} \\


\subsection{Записи}
\label{datatypes:record}

\textbf{Запись} это \emph{именованный кортеж}, имеющий именованные элементы,
которые называются \textbf{полями}. Тип записи определяется в виде атрибута 
модуля, например:

\begin{erlangru}
-record(Запись, {Поле1 [= Значение1],
                 ...
                 ПолеN [= ЗначениеN]}).
\end{erlangru}

Здесь имя записи \texttt{Запись} и имена полей \texttt{ПолеХ} --- атомы и каждое
\texttt{ПолеX} может получить необязательное значение по умолчанию
\texttt{ЗначениеX}. Это определение может быть помещено в любом месте модуля 
между определениями функций, но обязательно до первого использования. Если тип
записи используется в нескольких модулях, рекомендуется поместить его в
отдельный файл для включения.

Новая запись с типом \texttt{Запись} создаётся с помощью следующего выражения:

\begin{erlangru}
#Запись{Поле1=Выражение1, ..., ПолеK=ВыражениеK [, _=ВыражениеL]}
\end{erlangru}

Поля не обязательно должны идти в том же порядке, как и в определении записи.
Пропущенные поля получат свои соответствующие значения по умолчанию. Если 
использована последняя запись (подчёркивание равно \texttt{ВыражениеL}), то все
оставшиеся поля получат значение \texttt{ВыражениеL}. Поля без значений по
умолчанию и пропущенные поля получат значение \texttt{undefined}.

Значение поля можно получить используя выражение 
\texttt{Переменная\#Запись.Поле}.

\begin{erlang}
-module(employee).
-export([new/2]).
-record(person, {name, age, employed=erixon}).

new(Name, Age) -> #person{name=Name, age=Age}.
\end{erlang}

Здесь функция \texttt{employee:new/2} может быть использована в другом модуле,
который также должен включить файл, содержащий определение использованной здесь
записи \texttt{person}.

\texttt{\{P = employee:new(ernie,44)\}} \resultingin \texttt{\{person, ernie, 
	44, erixon\}} \\
\texttt{P\#person.age} \resultingin \texttt{44} \\
\texttt{P\#person.employed} \resultingin \texttt{erixon}

При работе с записями в интерпретаторе Erlang, можно использовать 
функции\linebreak
\texttt{rd(ИмяЗаписи, ОпределениеЗаписи)} и \texttt{rr(Модуль)} для того, чтобы
определить или загрузить новые определения записей. Подробнее читайте в
документации\linebreak\emph{Erlang Reference Manual}.



\subsection{Списки}
\label{datatypes:list}

\textbf{Список} --- это составной тип данных, который содержит \emph{переменное}
количество \textbf{термов}, заключённое в квадратных скобках.

\begin{erlangru}
[Терм1,...,ТермN]
\end{erlangru}

Каждый \texttt{ТермX} в списке называется \textbf{элементом}.  Количество
элементов в списке называется \textbf{длиной списка}. Как это принято в
функциональном программировании, первый элемент называется \textbf{головой}
списка, и остаток (начиная с 2\textsuperscript{го} элемента и до конца)
называется \textbf{хвостом} списка. Заметьте, что отдельные элементы в списке
не должны быть одинакового типа, хотя часто практикуется (и это, наверное, даже
удобно), иметь в списке элементы одинакового типа --- когда приходится работать
с элементами разных типов, обычно используются \textbf{записи}.

\begin{center}
\begin{tabular}{|>{\raggedright}p{124pt}|>{\raggedright}p{290pt}|}
\hline
\multicolumn{2}{|p{321pt}|}{Встроенные функции для работы со списками}  
\tabularnewline
\hline
\texttt{length(Список)} &
Возвращает длину \texttt{Списка} \tabularnewline
\hline
\texttt{hd(Список)} &
Возвращает 1\textsuperscript{й} элемент \texttt{Списка} (голову) \tabularnewline
\hline
\texttt{tl(Список)} &
Возвращает \texttt{Список} без 1\textsuperscript{го} элемента (хвост) \tabularnewline
\hline
\end{tabular}
\end{center}

Оператор "<вертикальная черта"> (\textbar{}, также в некоторых книгах по ФП он 
называется \texttt{cons}), отделяет ведущие элементы списка (один или более)
от остальных элементов. Например:

\texttt{[H | T]  = [1, 2, 3, 4, 5]} \resultingin \texttt{H=1} и 
	\texttt{T=[2, 3, 4, 5]} \\
\texttt{[X, Y | Z] = [a, b, c, d, e]} \resultingin \texttt{X=a},
	\texttt{Y=b} и \texttt{Z=[c, d, e]}

Список --- рекурсивная структура. Неявно список завершается ссылкой на пустой 
список, то есть~\texttt{[a,b]} это то же самое, как и \texttt{[a,b|[]]}. 
Список, выглядящий как \texttt{[a,b|c]} называется \textbf{плохо 
сформированным} (badly formed) и такой записи следует избегать (потому что
атом '\texttt{c}', завершающий структуру списка сам \emph{не} является списком).
Списки натурально способствуют рекурсивному функциональному программированию.

Например, следующая функция \texttt{sum} вычисляет сумму списка и функция 
\texttt{double} умножает каждый элемент списка на 2, при этом конструируя и 
возвращая новый список по ходу выполнения.

\begin{erlang}
sum([]) -> 0;
sum([H | T]) -> H + sum(T).

double([]) -> [];
double([H | T]) -> [H*2 | double(T)].
\end{erlang}

Определения функций выше представляют \emph{сопоставление с образцом}, которое
описано далее в главе \ref{patterns}. Образцы в такой записи часто встречаются 
в рекурсивном программировании, неявно предоставляя "<базовый случай"> (для
пустого списка в этих примерах).

Для работы со списками, оператор \texttt{++} соединяет вместе два списка 
(присоединяет второй аргумент к первому) и возвращает список-результат.  
Оператор  \texttt{--} создаёт список, который является копией первого
аргумента, за тем исключением, что для каждого элемента во втором списке его
первое вхождение в результат (если такое было) удаляется.

\texttt{[1,2,3] ++ [4,5]} \resultingin \texttt{[1,2,3,4,5]}
\texttt{[1,2,3,2,1,2] -- [2,1,2]} \resultingin \texttt{[3,1,2]}

Подборка функций, работающих со списками может быть найдена в модуле стандартной
библиотеки с названием \texttt{lists}.


\subsection{Строки}
\label{datatypes:string}

\textbf{Строки} --- это цепочки символов, заключённые между двойными кавычками,
на самом деле они хранятся в памяти, как списки целых чисел --- символов.

\texttt{"abcdefghi"} то же самое, как и 
	\texttt{[97,98,99,100,101,102,103,104,105]}

\texttt{""} то же самое, как и \texttt{[]}

Две строки, записанные подряд без разделительных знаков и операторов будут 
соединены в одну во время компиляции и не принесут дополнительных затрат по
соединению во время исполнения.

\texttt{"string" "42"} \resultingin \texttt{"string42"}


\subsection{Двоичные данные}
\label{datatypes:binary}

\textbf{Двоичные данные} --- это блок нетипизированной памяти, по умолчанию 
двоичные данные являются последовательностью 8-битных элементов (байтов).

\begin{erlangru}
<<Элемент1,...,ЭлементN>>
\end{erlangru}

Каждый \texttt{ЭлементX} указывается в виде

\begin{erlangru}
Значение[:Размер][/СписокСпецификацийТипа]
\end{erlangru}

\begin{center}
\begin{tabular}{|>{\raggedright}p{120pt}|>{\raggedright}p{120pt}|>{\raggedright}p{180pt}|}
\hline
\multicolumn{3}{|p{297pt}|}{Спецификация элемента двоичных данных} 
\tabularnewline
\hline
\texttt{Значение} &
\texttt{Размер} &
\texttt{СписокСпецификацийТипа}\tabularnewline
\hline
Выражение должно вычисляться в целое, действительное число или двоичные 
данные. & 
Выражение должно вычисляться в целое число & 
Последовательность необязательных спецификаторов типа, в любом порядке,
разделённых дефисами (-)\tabularnewline
\hline
\end{tabular}
\end{center}

% page break at this point, so turned into two chunks.
\begin{center}
\begin{tabular}{|>{\raggedright}p{90pt}|>{\raggedright}p{140pt}|>{\raggedright}p{200pt}|}
\hline
\multicolumn{3}{|p{297pt}|}{Спецификаторы типов}\tabularnewline
\hline
Тип данных &
\texttt{integer} \textbar{} \texttt{float} \textbar{} \texttt{binary} &
По умолчанию \texttt{integer}\tabularnewline
\hline
Наличие знака & 
\texttt{signed} \textbar{} \texttt{unsigned} & 
По умолчанию без знака, \texttt{unsigned}\tabularnewline
\hline
Порядок байтов (endianness) &
\texttt{big} \textbar{} \texttt{little} \textbar{} \texttt{native} &
Зависит от архитектуры процессора. По умолчанию \texttt{big}\tabularnewline
\hline
Единица измерения (группа битов) & 
\texttt{unit:}\emph{ЦелоеЧисло} & 
Разрешены значения от 1 до 256.
По умолчанию 1 для целых и действительных чисел и 8 для двоичных данных 
\tabularnewline
\hline
\end{tabular}
\end{center}

Значение \texttt{Размера} умножается на единицу измерения и даёт число бит, 
которые может занять данный сегмент двоичных данных. Каждый сегмент может
состоять из нуля или более битов, но общая сумма битов должна делиться нацело
на 8, иначе во время  исполнения возникнет ошибка \texttt{badarg}. Также
сегмент с типом \texttt{binary} должен иметь размер, делящийся нацело на 8.

Двоичные данные не могут вкладываться друг в друга.

\begin{erlang}
<<1, 17, 42>>       % <<1, 17, 42>>
<<"abc">>           % <<97, 98, 99>> (То же, что и <<$a, $b, $c>>)
<<1, 17, 42:16>>    % <<1,17,0,42>>
<<>>                % <<>>
<<15:8/unit:10>>    % <<0,0,0,0,0,0,0,0,0,15>>
<<(-1)/unsigned>>   % <<255>>
\end{erlang}


\section{Escape-последовательности}
\label{datatypes:escapeseq}

Escape-последовательности разрешено использовать в строках и атомах, которые 
заключены в кавычки, они позволяют ввести в исходный текст программы символ,
который другим способом ввести трудно или невозможно.

\begin{center}
\begin{tabular}{|>{\raggedright}p{90pt}|>{\raggedright}p{280pt}|}
\hline
\multicolumn{2}{|p{321pt}|}{Escape-последовательности}\tabularnewline
\hline
\textbackslash{}b & Backspace (удаление слева)\tabularnewline
\hline
\textbackslash{}d & Delete (удаление справа)\tabularnewline
\hline
\textbackslash{}e & Escape\tabularnewline
\hline
\textbackslash{}f & Новая страница (при печати)\tabularnewline
\hline
\textbackslash{}n & Новая строка\tabularnewline
\hline
\textbackslash{}r & Возврат курсора в начало\tabularnewline
\hline
\textbackslash{}s & Пробел\tabularnewline
\hline
\textbackslash{}t & Табуляция\tabularnewline
\hline
\textbackslash{}v & Вертикальная табуляция\tabularnewline
\hline
\textbackslash{}XYZ, \textbackslash{}XY, \textbackslash{}X &
Символ, записанный в восьмеричном представлении, как XYZ, XY или X\tabularnewline
\hline
\textbackslash{}\textasciicircum{}A .. \textbackslash{}\textasciicircum{}Z & 
Комбинации клавиш от Ctrl-A до Ctrl-Z\tabularnewline
\hline
\textbackslash{}\textasciicircum{}a .. \textbackslash{}\textasciicircum{}z &
Комбинации клавиш от Ctrl-A до Ctrl-Z\tabularnewline
\hline
\textbackslash{}' & Единичная кавычка \tabularnewline
\hline
\textbackslash{}\textbf{\texttt{"}} & Двойная кавычка \tabularnewline
\hline
\textbackslash{}\textbackslash{} & Обратная косая черта (\textbackslash) 
\tabularnewline
\hline
\end{tabular}
\end{center}


\section{Преобразования типов}

Erlang --- строго типизированный язык, то есть неявные автоматические 
преобразования типов в нём не происходят. Но есть ряд встроенных в стандартную
библиотеку функций, предназначенных для преобразования между типами данных
при участии программиста:

\begin{center}
\begin{tabular}{|>{\raggedright}p{63pt}|>{\raggedright}p{30pt}|>{\raggedright}p{35pt}|>{\raggedright}p{21pt}|>{\raggedright}p{21pt}|>{\raggedright}p{21pt}|>{\raggedright}p{30pt}|>{\raggedright}p{21pt}|>{\raggedright}p{35pt}|}
\hline
\multicolumn{9}{|p{243pt}|}{Преобразования типов}\tabularnewline
\hline
 & atom & integer & float & pid & fun & tuple & list & binary\tabularnewline
\hline
atom &  & - & - & - & - & - & X & X\tabularnewline
\hline
integer & - &  & X & - & - & - & X & X\tabularnewline
\hline
float & - & X &  & - & - & - & X & X\tabularnewline
\hline
pid & - & - & - &  & - & - & X & X\tabularnewline
\hline
fun & - & - & - & - &  & - & X & X\tabularnewline
\hline
tuple & - & - & - & - & - &  & X & X\tabularnewline
\hline
list & X & X & X & X & X & X &  & X\tabularnewline
\hline
binary & X & X & X & X & X & X & X & \tabularnewline
\hline
\end{tabular}
\end{center}

Встроенная функция \texttt{float/1} переводит целые числа в числа с 
плавающей точкой. Встроенные функции \texttt{round/1} и \texttt{trunc/1}
переводят числа с плавающей точкой обратно в целые, округляя или отбрасывая
дробную часть.

Функции \texttt{Тип\_to\_list/1} и \texttt{list\_to\_Тип/1}
переводят различные типы в списки (строки) и из списков.

Функции \texttt{term\_to\_binary/1} и \texttt{binary\_to\_term/1} переводят 
любое значение в закодированную двоичную форму и обратно (Подробнее: 
\url{http://erlang.org/doc/apps/erts/erl_ext_dist.html}).

\begin{erlang}
atom_to_list(hello)        % "hello"
list_to_atom("hello")      % hello
float_to_list(7.0)         % "7.00000000000000000000e+00"
list_to_float("7.000e+00") % 7.00000
integer_to_list(77)        % "77"
list_to_integer("77")      % 77
tuple_to_list({a, b ,c})   % [a,b,c]
list_to_tuple([a, b, c])   % {a,b,c}
pid_to_list(self())        % "<0.25.0>"
term_to_binary(<<17>>)     % <<131,109,0,0,0,1,17>>
term_to_binary({a, b ,c})  % <<131,104,3,100,0,1,97,100,0,1,98,
                           %   100,0,1,99>>
binary_to_term(<<131,104,3,100,0,1,97,100,0,1,98,100,0,1,99>>)  
                           % {a,b,c}
term_to_binary(math:pi())  % <<131,99,51,46,49,52,49,53,57,50,54,
                           % 53,51,...>>
\end{erlang}


\chapter{Pattern Matching}
\label{patterns}

\section{Variables}
\label{patterns:variables}

\textbf{Variables} are introduced as arguments to a function or as a
result of pattern matching. Variables begin with an uppercase letter
or underscore (\texttt{\_}) and may contain alphanumeric characters,
underscores and at-signs (\texttt{@}). Variables can only be bound
(assigned) once.

\begin{erlang}
Abc
A_long_variable_name
AnObjectOrientatedVariableName
_Height
\end{erlang}

An \textbf{anonymous variable} is denoted by a single underscore (\texttt{\_})
and can be used when a variable is required but its value can be
ignored.

\begin{erlang}
[H|_] = [1,2,3]         % H=1 and the rest is ignored
\end{erlang}

Variables beginning with underscore like \texttt{\_Height} are normal
variables, not anonymous. They are however ignored by the compiler in
the sense that they will not generate any warnings for unused
variables. Thus it is possible to write:

\begin{erlang}
member(_Elem, []) ->
    false.
\end{erlang}

instead of:

\begin{erlang}
member(_, []) ->
    false.
\end{erlang}

which can make for more readable code.

The \textit{scope} for a variable is its function clause. Variables
bound in a branch of an \texttt{if}, \texttt{case}, or
\texttt{receive} expression must be bound in all branches to have a
value outside the expression, otherwise they will be regarded as
\textit{unsafe} (possibly undefined) outside the expression.


\section{Pattern Matching}
A \textbf{pattern} has the same structure as a term but may contain
new unbound variables.

\begin{erlang}
Name1
[H|T]
{error,Reason}
\end{erlang}

Patterns occur in \textit{function heads}, \textit{case},
\textit{receive}, and \textit{try} expressions and in match operator
(\texttt{=}) expressions. Patterns are evaluated through
\textbf{pattern matching} against an expression and this is how
variables are defined and bound.

\begin{erlang}
Pattern = Expr
\end{erlang}

Both sides of the expression must have the same structure. If the
matching succeeds, all unbound variables, if any, in the pattern
become bound. If the matching fails, a \texttt{badmatch} run-time
error will occur.


\begin{minted}{console}
> {A, B} = {answer, 42}.
{answer,42}
> A.
answer
> B.
42
\end{minted}

\subsection{Match operator (\texttt{=}) in patterns}
If \texttt{Pattern1} and \texttt{Pattern2} are valid patterns, then
the following is also a valid pattern:

\begin{erlang}
Pattern1 = Pattern2
\end{erlang}

The \texttt{=} introduces an \textbf{alias} which when matched against an
expression, both \texttt{Pattern1} and \texttt{Pattern2} are matched
against it. The purpose of this is to avoid the reconstruction of terms.

\begin{erlang}
foo({connect,From,To,Number,Options}, To) ->
    Signal = {connect,From,To,Number,Options},
    fox(Signal),
    ...;
\end{erlang}

which can be written more efficiently as:

\begin{erlang}
foo({connect,From,To,Number,Options} = Signal, To) ->
    fox(Signal),
    ...;
\end{erlang}


\subsection{String prefix in patterns}
When matching strings, the following is a valid pattern:

\begin{erlang}
f("prefix" ++ Str) -> ...
\end{erlang}

which is equivalent to and easier to read than:

\begin{erlang}
f([$p,$r,$e,$f,$i,$x | Str]) -> ...
\end{erlang}

You can only use strings as prefix expressions; patterns such as \texttt{Str ++ "postfix"} are not allowed. 

\subsection{Expressions in patterns}
An arithmetic expression can be used within a pattern, provided it
only uses numeric or bitwise operators and its value can be evaluated
to a constant at compile-time.

\begin{erlang}
case {Value, Result} of
    {?Threshold+1, ok} -> ...   % ?Threshold is a macro
\end{erlang}

\subsection{Matching binaries}

\begin{erlang}
Bin = <<1, 2, 3>>               % <<1,2,3>> All elements are 8-bit bytes
<<A, B, C>> = Bin               % A=1, B=2 and C=3
<<D:16, E>> = Bin               % D=258 and E=3
<<F, G/binary>> = Bin           % F=1 and G=<<2,3>>
\end{erlang}

In the last line, the variable \texttt{G} of unspecified size matches the
rest of the binary \texttt{Bin}.

Always put a space between (\texttt{=}) and (\verb|<<|) so as to
avoid confusion with the (\texttt{=<}) operator.


\chapter{Functions}

\section{Function definition}
A function is defined as a sequence of one or more \textbf{function
clauses}. The function name is an atom.

\vspace*{4pt}
\begin{erlang}
Func(Pattern11,...,Pattern1N) [when GuardSeq1] -> Body1;
    ...;
    ...;
Func(PatternK1,...,PatternKN) [when GuardSeqK] -> BodyK.

\end{erlang}
\vspace*{4pt}

The function clauses are separated by semicolons (\texttt{;}) and
terminated by full stop (\texttt{.}). A function clause consists of a
\textbf{clause head} and a \textbf{clause body} separated by an arrow
(\texttt{->}). A clause head consists of the function name (an atom),
arguments within parentheses and an optional guard sequence beginning
with the keyword \texttt{when}.  Each argument is a pattern.  A clause
body consists of a sequence of expressions separated by commas
(\texttt{,}).

\vspace*{4pt}
\begin{erlang}
Expr1,
...,
ExprM
\end{erlang}
\vspace*{4pt}

The number of arguments \texttt{N} is the \textbf{arity} of the
function. A function is uniquely defined by the module name, function name
and arity. Two different functions in the same module with different
arities may have the same name. A function \texttt{Func} in \texttt{Module}
with arity \texttt{N} is often denoted as \texttt{Module:Func/N}.

\vspace*{4pt}
\begin{erlang}
-module(mathStuff).
-export([area/1]).

area({square, Side}) -> Side * Side;
area({circle, Radius}) -> math:pi() * Radius * Radius;
area({triangle, A, B, C}) ->
    S = (A + B + C)/2,
    math:sqrt(S*(S-A)*(S-B)*(S-C)).
\end{erlang}

% push this section onto the next page (orphan line)
\newpage
\section{Function calls}
A function is called using:

\begin{erlang}
[Module:]Func(Expr1, ..., ExprN)
\end{erlang}

\texttt{Module} evaluates to a module name and \texttt{Func} to
a function name or a \textit{fun}. When calling a function in another
module, the module name must be provided and the function must be
exported. This is referred to as a \textbf{fully qualified function name}.

\begin{erlang}
lists:keysearch(Name, 1, List)
\end{erlang}

The module name can be omitted if \texttt{Func} evaluates to the name
of a local function, an imported function, or an auto-imported
BIF.  In such cases, the function is called using an \textbf{implicitly qualified function name}.

Before calling a function the arguments \texttt{ExprX} are
evaluated.  If the function cannot be found, an \texttt{undef} run-time
error will occur. Next the function clauses are scanned sequentially
until a clause is found such that the patterns in the clause head can
be successfully matched against the given arguments and that the guard
sequence, if any, is true. If no such clause can be found, a
\texttt{function\_clause} run-time error will occur.

If a matching clause is found, the corresponding clause body is evaluated,
i.e.~the expressions in the body are evaluated sequentially and the
value of the last expression is returned.

The fully qualified function name must be used when calling a function
with the same name as a BIF (built-in function, see section
\ref{functions:bifs}). The compiler does not allow defining a function
with the same name as an imported function. When calling a local
function, there is a difference between using the implicitly or fully
qualified function name, as the latter always refers to the latest
version of the module (see chapter \ref{code}).


\section{Expressions}
\label{functions:expressions}
An \textbf{expression} is either a term or the invocation of an
operator, for example:

\begin{erlang}
Term
op Expr
Expr1 op Expr2
(Expr)
begin
   Expr1,
   ...,
   ExprM            % no comma (,) before end
end

\end{erlang}

There are both \textit{unary} and \textit{binary} operators. The
simplest form of expression is a term, i.e.~an \textit{integer},
\textit{float}, \textit{atom}, \textit{string}, \textit{list} or
\textit{tuple} and the return value is the term itself. An expression
may contain \textit{macro} or \textit{record} expressions which will
expanded at compile time.

Parenthesised expressions are useful to override operator precedence (see section \ref{functions:expressions:precedence}):

\begin{erlang}
1 + 2 * 3           % 7
(1 + 2) * 3         % 9
\end{erlang}

Block expressions within \texttt{begin...end} can be used to group a
sequence of expressions and the return value is the value of the last
expression \texttt{ExprM}.

All subexpressions are evaluated before the expression itself is
evaluated, but the order in which subexpressions are evaluated is undefined.

Most operators can only be applied to arguments of a certain type. For
example, arithmetic operators can only be applied to integers or
floats. An argument of the wrong type will cause a \texttt{badarg}
run-time error.


\subsection{Term comparisons}
\begin{erlang}
Expr1 op Expr2
\end{erlang}

A \textbf{term comparison} returns a \textit{boolean} value,
in the form of atoms \texttt{true} or \texttt{false}.

\begin{center}
\begin{tabular}{|>{\raggedright}p{40pt}|>{\raggedright}p{105pt}|>{\raggedright}p{26pt}|>{\raggedright}p{124pt}|}
\hline
\multicolumn{4}{|p{297pt}|}{Comparison operators}\tabularnewline
\hline
\texttt{==} & Equal to & \texttt{=<} & Less than or equal to\tabularnewline
\hline
\texttt{/=} & Not equal to & \texttt{<} & Less than\tabularnewline
\hline
\texttt{=:=} & Exactly equal to & \texttt{>}= & Greater than or equal to\tabularnewline
\hline
\texttt{=/=} & Exactly not equal to & \texttt{>} & Greater than\tabularnewline
\hline
\end{tabular}
\end{center}

\begin{erlang}
1==1.0              % true
1=:=1.0             % false
1 > a               % false
\end{erlang}

The arguments may be of different data types. The following order is
defined:

\texttt{number < atom < reference < fun < port < pid < tuple < list < binary}

Lists are compared element by element. Tuples are ordered by size, two
tuples with the same size are compared element by element. When
comparing an integer and a float, the integer is first converted to a
float. In the case of \texttt{=:=} or \texttt{=/=} there is no type conversion.


\subsection{Arithmetic expressions}

\begin{erlang}
op Expr
Expr1 op Expr2
\end{erlang}

An \textbf{arithmetic expression} returns the result after applying
the operator.

\begin{center}
\begin{tabular}{|>{\raggedright}p{35pt}|>{\raggedright}p{145pt}|>{\raggedright}p{128pt}|}
\hline
\multicolumn{3}{|p{309pt}|}{Arithmetic operators}\tabularnewline
\hline
\texttt{+}  & Unary +  & \texttt{Integer \textbar{} Float} \tabularnewline
\hline
\texttt{-}  & Unary -  & \texttt{Integer \textbar{} Float}\tabularnewline
\hline
\texttt{+}  & Addition & \texttt{Integer} \textbar{} \texttt{Float}\tabularnewline
\hline
\texttt{-}  & Subtraction & \texttt{Integer} \textbar{} \texttt{Float}\tabularnewline
\hline
\texttt{*}  & Multiplication & \texttt{Integer} \textbar{} \texttt{Float}\tabularnewline
\hline
\texttt{/}  & Floating point division  & \texttt{Integer} \textbar{} \texttt{Float}\tabularnewline
\hline
\texttt{bnot}  & Unary bitwise not  & \texttt{Integer} \tabularnewline
\hline
\texttt{div}  & Integer division  & \texttt{Integer}\tabularnewline
\hline
\texttt{rem}  & Integer remainder of X/Y  & \texttt{Integer} \tabularnewline
\hline
\texttt{band}  & Bitwise and  & \texttt{Integer}\tabularnewline
\hline
\texttt{bor}  & Bitwise or  & \texttt{Integer} \tabularnewline
\hline
\texttt{bxor}  & Arithmetic bitwise xor  & \texttt{Integer}\tabularnewline
\hline
\texttt{bsl}  & Arithmetic bitshift left  & \texttt{Integer} \tabularnewline
\hline
\texttt{bsr}  & Bitshift right  & \texttt{Integer}\tabularnewline
\hline
\end{tabular}
\end{center}

\begin{erlang}
+1                  % 1
4/2                 % 2.00000
5 div 2             % 2
5 rem 2             % 1
2#10 band 2#01      % 0
2#10 bor 2#01       % 3
\end{erlang}


\subsection{Boolean expressions}

\begin{erlang}
op Expr
Expr1 op Expr2
\end{erlang}

A \textbf{boolean expression} returns the value \texttt{true} or
\texttt{false} after applying the operator.

\begin{center}
\begin{tabular}{|>{\raggedright}p{79pt}|>{\raggedright}p{241pt}|}
\hline
\multicolumn{2}{|p{321pt}|}{Boolean operators}\tabularnewline
\hline
\texttt{not}  & Unary logical not \tabularnewline
\hline
\texttt{and}  & Logical and \tabularnewline
\hline
\texttt{or}  & Logical or \tabularnewline
\hline
\texttt{xor}  & Logical exclusive or\tabularnewline
\hline
\end{tabular}
\end{center}

\begin{erlang}
not true            % false
true and false      % false
true xor false      % true
\end{erlang}


\subsection{Short-circuit boolean expressions}

\begin{erlang}
Expr1 orelse Expr2
Expr1 andalso Expr2
\end{erlang}

These are boolean expressions where \texttt{Expr2} is evaluated only
if necessary. In an \texttt{orelse} expression \texttt{Expr2} will be
evaluated only if \texttt{Expr1} evaluates to false. In an
\texttt{andalso} expression \texttt{Expr2} will be evaluated only if
\texttt{Expr1} evaluates to true.

\begin{erlang}
if A >= 0 andalso math:sqrt(A) > B -> ...

if is_list(L) andalso length(L) == 1 -> ...
\end{erlang}


\subsection{Operator precedences}
\label{functions:expressions:precedence}
In an expression consisting of subexpressions the operators will be
applied according to a defined \textbf{operator precedence} order:

\begin{center}
\begin{tabular}{|>{\raggedright}p{221pt}|>{\raggedright}p{99pt}|}
\hline
\multicolumn{2}{|p{321pt}|}{Operator precedence (from high to low)}\tabularnewline
\hline
\texttt{:} ~ &  \tabularnewline
\hline
\texttt{\#} ~ &  \tabularnewline
\hline
\texttt{Unary + - bnot not ~} &  \tabularnewline
\hline
\texttt{/ * div rem band and}  & Left associative \tabularnewline
\hline
\texttt{+ - bor bxor bsl bsr or xor} & Left associative \tabularnewline
\hline
\texttt{++ --}  & Right associative \tabularnewline
\hline
\texttt{== /= =< < >= > =:= =/=} & \tabularnewline
\hline
\texttt{andalso}  &  \tabularnewline
\hline
\texttt{orelse} &  \tabularnewline
\hline
\texttt{= !}  & Right associative \tabularnewline
\hline
\texttt{catch ~} &  \tabularnewline
\hline
\end{tabular}
\end{center}

The operator with the highest priority is evaluated first. Operators
with the same priority are evaluated according to their
\textbf{associativity}.  The left associative arithmetic operators
are evaluated left to right:

\texttt{6 + 5 * 4 - 3 / 2  \resultingin  6 + 20 - 1.5  \resultingin  26 - 1.5  \resultingin  24.5}


\section{Compound expressions}


\subsection{If}

\begin{erlang}
if
    GuardSeq1 ->
        Body1;
    ...;
    GuardSeqN ->
        BodyN                   % Note no semicolon (;) before end
end
\end{erlang}

The branches of an \texttt{if} expression are scanned sequentially
until a guard sequence \texttt{GuardSeq} which evaluates to
\texttt{true} is found.  The corresponding \texttt{Body} (sequence
of expressions separated by commas) is then evaluated.  The return value of
\texttt{Body} is the return value of the \texttt{if} expression.

If no guard sequence is true, an \texttt{if\_clause} run-time error
will occur. If necessary, the guard expression \texttt{true} can be used in the
last branch, as that guard sequence is always true (known as a ``catch
all'').

\begin{erlang}
is_greater_than(X, Y) ->
    if
        X>Y ->
            true;
        true ->                 % works as an 'else' branch
            false
    end
\end{erlang}

It should be noted that pattern matching in function clauses can be used to replace \texttt{if} cases (most of the time).
Overuse of \texttt{if} sentences withing function bodies is considered a bad Erlang practice.

\subsection{Case}

Case expressions provide for inline pattern matching, similar to the way in which function clauses are matched.

\begin{erlang}
case Expr of
    Pattern1 [when GuardSeq1] ->
        Body1;
        ...;
    PatternN [when GuardSeqN] ->
        BodyN                   % Note no semicolon (;) before end
end
\end{erlang}

The expression \texttt{Expr} is evaluated and the patterns
\texttt{Pattern1}...\texttt{PatternN} are sequentially matched against the result. If a
match succeeds and the optional guard sequence \texttt{GuardSeqX} is
\texttt{true}, then the corresponding \texttt{BodyX} is evaluated. The return value
of \texttt{BodyX} is the return value of the case expression.

If there is no matching pattern with a true guard sequence, a
\texttt{case\_clause} run-time error will occur.

\begin{erlang}
is_valid_signal(Signal) ->
    case Signal of
        {signal, _What, _From, _To} ->
            true;
        {signal, _What, _To} ->
            true;
        _Else ->                % 'catch all'
            false
    end.
\end{erlang}


\subsection{List comprehensions}
List comprehensions are analogous to the \texttt{setof} and
\texttt{findall} predicates in Prolog.

\begin{erlang}
[Expr || Qualifier1,...,QualifierN]
\end{erlang}

\texttt{Expr} is an arbitrary expression, and each \texttt{QualifierX}
is either a \textbf{generator} or a \textbf{filter}. A generator is
written as:

\begin{erlang}
Pattern <- ListExpr
\end{erlang}

where \texttt{ListExpr} must be an expression which evaluates to a
list of terms. A filter is an expression which evaluates to
\texttt{true} or \texttt{false}. Variables in list generator
expressions \textit{shadow} variables in the function clause
surrounding the list comprehension.

The qualifiers are evaluated from left to right, the generators
creating values and the filters constraining them. The list
comprehension then returns a list where the elements are the result of
evaluating \texttt{Expr} for each combination of the resulting values.

\begin{minted}{console}
> [{X, Y} || X <- [1,2,3,4,5,6], X > 4, Y <- [a,b,c]].
[{5,a},{5,b},{5,c},{6,a},{6,b},{6,c}]
\end{minted}


\section{Guard sequences}
A \textbf{guard sequence} is a set of \textbf{guards} separated by
semicolons (\texttt{;}). The guard sequence is \texttt{true} if at
least one of the guards is \texttt{true}.

\begin{erlang}
Guard1; ...; GuardK
\end{erlang}

A \textbf{guard} is a set of \textbf{guard expressions} separated by
commas (\texttt{,}). The guard is \texttt{true} if all guard
expressions evaluate to \texttt{true}.

\begin{erlang}
GuardExpr1, ..., GuardExprN
\end{erlang}

The permitted \textbf{guard expressions} (sometimes called guard
tests) are a subset of valid Erlang expressions, since the
evaluation of a guard expression must be guaranteed to be free of side-effects.

\begin{center}
\begin{tabular}{|>{\raggedright}p{154pt}|>{\raggedright}p{166pt}|}
\hline
\multicolumn{2}{|p{321pt}|}{Valid guard expressions:}\tabularnewline
\hline
\multicolumn{2}{|p{321pt}|}{The atom \texttt{true};}\tabularnewline
\hline
\multicolumn{2}{|p{321pt}|}{Other constants (terms and bound variables), are all regarded
as \texttt{false};}\tabularnewline
\hline
\multicolumn{2}{|p{321pt}|}{Term comparisons;}\tabularnewline
\hline
\multicolumn{2}{|p{321pt}|}{Arithmetic and boolean expressions;}\tabularnewline
\hline
\multicolumn{2}{|p{321pt}|}{Calls to the BIFs specified below.}\tabularnewline
\hline
Type test BIFs & Other BIFs allowed in guards:\tabularnewline
\hline
\texttt{is\_atom/1} & \texttt{abs(Integer} \textbar{} \texttt{Float)}\tabularnewline
\hline
\texttt{is\_constant/1} & \texttt{float(Term)}\tabularnewline
\hline
\texttt{is\_integer/1} & \texttt{trunc(Integer} \textbar{} \texttt{Float)}\tabularnewline
\hline
\texttt{is\_float/1} & \texttt{round(Integer} \textbar{} \texttt{Float)}\tabularnewline
\hline
\texttt{is\_number/1} & \texttt{size(Tuple} \textbar{} \texttt{Binary)}\tabularnewline
\hline
\texttt{is\_reference/1} & \texttt{element(N, Tuple)}\tabularnewline
\hline
\texttt{is\_port/1} & \texttt{hd(List)}\tabularnewline
\hline
\texttt{is\_pid/1} & \texttt{tl(List)}\tabularnewline
\hline
\texttt{is\_function/1} & \texttt{length(List)}\tabularnewline
\hline
\texttt{is\_tuple/1} & \texttt{self()}\tabularnewline
\hline
\texttt{is\_record/2} The 2\textsuperscript{nd} argument is \linebreak{}
the record name & \texttt{node(})\tabularnewline
\hline
\texttt{is\_list/1} & \texttt{node(Pid} \textbar{} \texttt{Ref} \textbar \texttt{Port)}\tabularnewline
\hline
\texttt{is\_binary/1} & \tabularnewline
\hline
\end{tabular}
\end{center}

A small example:

\begin{erlang}
fact(N) when N>0 ->             % first clause head
    N * fact(N-1);              % first clause body
fact(0) ->                      % second clause head
    1.                          % second clause body
\end{erlang}


\section{Tail recursion}
If the last expression of a function body is a function call, a
\textbf{tail recursive} call is performed in such a way that no system
resources (like the call stack) are consumed. This means that an
infinite loop like a server can be programmed provided it only uses
tail recursive calls.

The function \texttt{fact/1} above could be rewritten using tail
recursion in the following manner:

 \begin{erlang}
fact(N) when N>1 -> fact(N, N-1);
fact(N) when N==1; N==0 -> 1.

fact(F,0) -> F;                 % The variable F is used as an accumulator
fact(F,N) -> fact(F*N, N-1).
\end{erlang}


\section{Funs}
\label{functions:funs}
A \textbf{fun} defines a \textit{functional object}. Funs make it
possible to pass an entire function, not just the function name, as an
argument. A `fun' expression begins with the keyword \texttt{fun} and
ends with the keyword \texttt{end} instead of a full stop
(\texttt{.}).  Between these should be a regular function
declaration, except that no function name is specified.

\begin{erlang}
fun
    (Pattern11,...,Pattern1N) [when GuardSeq1] ->
        Body1;
        ...;
    (PatternK1,...,PatternKN) [when GuardSeqK] ->
        BodyK
end
\end{erlang}

Variables in a \texttt{fun} head \textit{shadow} variables in the function
clause surrounding the \texttt{fun} but variables bound in a \texttt{fun} body are local
to the body.  The return value of the expression is the resulting function. The expression
\texttt{fun Name/N is} equivalent to:

\begin{erlang}
fun (Arg1,...,ArgN) -> Name(Arg1,...,ArgN) end
\end{erlang}

The expression \texttt{fun Module:Func/Arity} is also allowed, provided that \texttt{Func} is exported
from \texttt{Module}.

\begin{erlang}
Fun1 = fun (X) -> X+1 end.
Fun1(2)         % 3

Fun2 = fun (X) when X>=1000 -> big; (X) -> small end.
Fun2(2000)      % big
\end{erlang}

Since a \texttt{fun} is anonymous, i.e.~there is no function name in the
definition of the \texttt{fun}, the definition of a recursive \texttt{fun} has to be
done in two steps.  This example shows how to define the function
\texttt{sum(List)} (see section \ref{datatypes:list}) as a \texttt{fun}.

\begin{erlang}
Sum1 = fun ([], _Foo) -> 0;([H|T], Foo) -> H + Foo(T, Foo) end.
Sum = fun (List) -> Sum1(List, Sum1) end.
Sum([1,2,3,4,5])    % 15
\end{erlang}

The definition of \texttt{Sum1} is done in a way such that it takes \textit{itself} as a parameter, matched to \texttt{\_Foo} (empty list) or \texttt{Foo},
which it then calls recursively.  The definition of \texttt{Sum} calls \texttt{Sum1}, also passing \texttt{Sum1} as a parameter.

\section{BIFs --- Built-in functions}
\label{functions:bifs}
The \textbf{built-in functions}, BIFs, are implemented in the C code of
the runtime system and do things that are difficult or impossible to
implement in Erlang. Most of the built-in functions belong to the
module \texttt{erlang} but there are also built-in functions that belong
to other modules like \texttt{lists} and \texttt{ets}. The most
commonly used BIFs belonging to the module \texttt{erlang} are
\textbf{auto-imported}, i.e.~they do not need to be prefixed with the
module name.

\begin{center}
\begin{tabular}{|>{\raggedright}p{103pt}|>{\raggedright}p{217pt}|}
\hline
\multicolumn{2}{|p{321pt}|}{Some useful BIFs}\tabularnewline
\hline
\texttt{date()} & Returns today's date as \texttt{\{Year, Month, Day\}}\tabularnewline
\hline
\texttt{now()} & Returns current time in microseconds. System dependent\tabularnewline
\hline
\texttt{time()} & Returns current time as \texttt{\{Hour, Minute, Second\}} System dependent\tabularnewline
\hline
\texttt{halt()} & Stops the Erlang system\tabularnewline
\hline
\texttt{processes()} & Returns a list of all processes currently known to the system\tabularnewline
\hline
\texttt{process\_info(Pid)} & Returns a dictionary containing information about \texttt{Pid}\tabularnewline
\hline
\texttt{Module:module\_info()} & Returns a dictionary containing information about the code
in Module\tabularnewline
\hline
\end{tabular}
\end{center}

A \textbf{dictionary} is a list of \texttt{\{Key, Value\} terms (see
also section \ref{processes:dicts}).}

\begin{erlang}
size({a, b, c})             % 3
atom_to_list('Erlang')      % "Erlang"
date()                      % {2013,5,27}
time()                      % {01,27,42}
\end{erlang}

\chapter{Processes}
\label{processes}

A \textbf{process} corresponds to one \textit{thread of control}.
Erlang permits very large numbers of concurrent processes,
each executing like it had an own virtual processor.  When a process
executing \texttt{functionA} calls another \texttt{functionB}, it
will wait until \texttt{functionB} is finished and then retrieve its
result. If instead it \textit{spawns} another process executing
\texttt{functionB}, both will continue in parallel
(concurrently). \texttt{functionA} will not wait for \texttt{functionB}
and the only way they can communicate is through \textit{message passing}.

Erlang processes are light-weight with a small memory footprint, fast to
create and shut-down, and the scheduling overhead is low.  A
\textbf{process identifier}, \texttt{Pid}, identifies a process. The
BIF \texttt{self/0} returns the \texttt{Pid} of the calling process.


\section{Process creation}
A process is created using the BIF \texttt{spawn/3}.

\begin{erlang}
spawn(Module, Func, [Expr1, ..., ExprN])
\end{erlang}

\texttt{Module} should evaluate to a module name and \texttt{Func} to
a function name in that module. The list \texttt{Expr1}$...$\texttt{ExprN} are the
arguments to the function.  \texttt{spawn} creates a new process and
returns the process identifier, \texttt{Pid}. The new process starts
by executing:

\begin{erlang}
Module:Func(Expr1, ..., ExprN)
\end{erlang}

The function \texttt{Func} has to be exported even if it is spawned by
another function in the same module. There are other spawn BIFs, for
example \texttt{spawn/4} for spawning a process on another node.


\section{Registered processes}
A process can be associated with a name. The name must be an atom and
is automatically unregistered if the process terminates. Only static
(cyclic) processes should be registered.

\begin{center}
\begin{tabular}{|>{\raggedright}p{117pt}|>{\raggedright}p{204pt}|}
\hline
\multicolumn{2}{|p{321pt}|}{Name registration BIFs}\tabularnewline
\hline
\texttt{register(Name, Pid)}  & Associates the atom \texttt{Name} with the process \texttt{Pid}\tabularnewline
\hline
\texttt{registered()}  & Returns a list of names which have been registered \tabularnewline
\hline
\texttt{whereis(Name)}  & Returns the \texttt{Pid} registered under \texttt{Name} or \texttt{undefined} if the name
is not registered\tabularnewline
\hline
\end{tabular}
\end{center}


\section{Process communication}
Processes communicate by sending and receiving
\textbf{messages}. Messages are sent using the send operator
(\texttt{!}) and are received using \texttt{receive}. Message passing
is asynchronous and reliable, i.e.~the message is guaranteed to
eventually reach the recipient, provided that the recipient exists.

\subsection{Send}
\begin{erlang}
Pid ! Expr
\end{erlang}

The send (\texttt{!}) operator sends the value of \texttt{Expr} as a
message to the process specified by \texttt{Pid} where it will be
placed last in its \textbf{message queue}.  The value of \texttt{Expr}
is also the return value of the (\texttt{!}) expression. \texttt{Pid}
must evaluate to a process identifier, a registered name or a tuple
\texttt{\{Name,Node\}}, where \texttt{Name} is a registered process at
\texttt{Node} (see chapter \ref{distribution}). The message sending operator
(\texttt{!}) never fails, even if it addresses a non-existent process.


\subsection{Receive}

\begin{erlang}
receive
    Pattern1 [when GuardSeq1] ->
        Body1;
    ...
    PatternN [when GuardSeqN] ->
        BodyN                   % Note no semicolon (;) before end
end
\end{erlang}

This expression receives messages sent to the process using the send
operator (\texttt{!}). The patterns \texttt{PatternX} are sequentially
matched against the first message in time order in the message queue,
then the second and so on.  If a match succeeds and the optional guard
sequence \texttt{GuardSeqX} is true, then the message is removed from
the message queue and the corresponding \texttt{BodyX} is
evaluated.  It is the order of the pattern clauses that decides the
order in which messages will be received prior to the order in which
they arrive.  This is called \textit{selective receive}. The
return value of \texttt{BodyX} is the return value of the receive
expression.

\texttt{receive} never fails. The process may be suspended, possibly
indefinitely, until a message arrives that matches one of the patterns
and with a true guard sequence.

\newpage
\begin{erlang}
wait_for_onhook() ->
    receive
        onhook ->
            disconnect(),
            idle();
        {connect, B} ->
            B ! {busy, self()},
            wait_for_onhook()
    end.
\end{erlang}


\subsection{Receive with timeout}

\begin{erlang}
receive
    Pattern1 [when GuardSeq1] ->
        Body1;
        ...;
    PatternN [when GuardSeqN] ->
        BodyN
after
    ExprT ->
        BodyT
end
\end{erlang}

\texttt{ExprT} should evaluate to an integer between \texttt{0} and
\texttt{16\#ffffffff} (the value must fit in 32 bits). If no matching
message has arrived within \texttt{ExprT} milliseconds, then
\texttt{BodyT} will be evaluated and its return value becomes the
return value of the receive expression.

\begin{erlang}
wait_for_onhook() ->
    receive
        onhook ->
            disconnect(),
            idle();
        {connect, B} ->
            B ! {busy, self()},
            wait_for_onhook()
    after
        60000 ->
            disconnect(),
            error()
    end.
\end{erlang}

A \texttt{receive...after} expression with no branches can be used to
implement simple timeouts.

\begin{erlang}
receive
after
    ExprT ->
        BodyT
end
\end{erlang}

\begin{center}
\begin{tabular}{|>{\raggedright}p{47pt}|>{\raggedright}p{273pt}|}
\hline
\multicolumn{2}{|p{321pt}|}{Two special cases for the timeout value \texttt{ExprT}}\tabularnewline
\hline
\texttt{infinity} & This is equivalent to not using a timeout and can be useful for timeout
values that are calculated at run-time\tabularnewline
\hline
\texttt{0} & If there is no matching message in the mailbox, the timeout will occur immediately\tabularnewline
\hline
\end{tabular}
\end{center}


\section{Process termination}
\label{processes:termination}
A process always terminates with an \textbf{exit reason} which may be
any term.  If a process terminates normally, i.e.~it has run
to the end of its code, then the reason is the atom \texttt{normal}. A process
can terminate itself by calling one of the following BIFs.

\begin{erlang}
exit(Reason)

erlang:error(Reason)

erlang:error(Reason, Args)
\end{erlang}

A process terminates with the exit reason \texttt{\{Reason,Stack\}} when a
run-time error occurs.

A process may also be terminated if it receives an exit signal with
a reason other than \texttt{normal} (see section \ref{processes:recvexitsignals}).


\section{Process links}
\label{processes:links}
Two processes can be \textbf{linked} to each other. Links are
bidirectional and there can only be one link between two distinct processes (unique \texttt{Pid}s). A
process with \texttt{Pid1} can link to a process with \texttt{Pid2}
using the BIF \texttt{link(Pid2)}.  The BIF \texttt{spawn\_link(Module, Func, Args)}
spawns and links a process in one atomic operation.

A link can be removed using the BIF \texttt{unlink(Pid)}.


\subsection{Error handling between processes}
When a process terminates it will send \textbf{exit signals} to all
processes that it is linked to.  These in turn will also be terminated
\textit{or handle the exit signal in some way}. This feature can be
used to build hierarchical program structures where some processes are
supervising other processes, for example restarting them if they
terminate abnormally.


\subsection{Sending exit signals}
\label{processes:sendexitsignals}
A process always terminates with an exit reason which is sent as an
exit signal to all linked processes. The BIF \texttt{exit(Pid, Reason)} sends
an exit signal with the reason \texttt{Reason} to \texttt{Pid}, without
affecting the calling process.


\subsection{Receiving exit signals}
\label{processes:recvexitsignals}
If a process receives an exit signal with an exit reason other than
\texttt{normal} it will also be terminated, and will send exit signals with the
same exit reason to its linked processes.  An exit signal with reason
\texttt{normal} is ignored.  This behaviour can be changed using the BIF
\texttt{process\_flag(trap\_exit, true)}.

The process is then able to \textbf{trap exits}.  This means that an
exit signal will be transformed into a message \texttt{\{'EXIT', FromPid, Reason\}} which is
put into the process's mailbox and can be handled by the process like a regular
message using \texttt{receive}.

However, a call to the BIF \texttt{exit(Pid, kill)} unconditionally
terminates the process \texttt{Pid} regardless whether it is able to
trap exit signals or not.


\section{Monitors}
A process \texttt{Pid1} can create a \textbf{monitor} for
\texttt{Pid2} using the BIF:

\begin{erlang}
erlang:monitor(process, Pid2)
\end{erlang}

which returns a reference \texttt{Ref}. If \texttt{Pid2} terminates
with exit reason \texttt{Reason}, a message as follows will be sent to
\texttt{Pid1}:

\begin{erlang}
{'DOWN', Ref, process, Pid2, Reason}
\end{erlang}

If \texttt{Pid2} does not exist, the \texttt{'DOWN'} message is sent
immediately with \texttt{Reason} set to \texttt{noproc}.  Monitors are
unidirectional in that if \texttt{Pid1} monitors \texttt{Pid2} then it
will receive a message when \texttt{Pid2} dies but \texttt{Pid2} will
\textbf{not} receive a message when \texttt{Pid1} dies. Repeated calls
to \texttt{erlang:monitor(process, Pid)} will create several,
independent monitors and each one will be sent a \texttt{'DOWN'} message
when \texttt{Pid} terminates.

A monitor can be removed by calling \texttt{erlang:demonitor(Ref)}. It
is possible to create monitors for processes with registered names,
also at other nodes.


\section{Process priorities}
The BIF \texttt{process\_flag(priority, Prio)} defines the priority of
the current process. \texttt{Prio} may have the value \texttt{normal},
which is the default, \texttt{low}, \texttt{high} or \texttt{max}. 

Modifying a process's priority is discouraged and should only be done in 
special circumstances.  A problem that requires changing process priorities
can generally be solved by another approach.


\section{Process dictionary}
\label{processes:dicts}
Each process has its own process dictionary which is a list of
\texttt{\{Key, Value\}} terms.

\begin{center}
\begin{tabular}{|>{\raggedright}p{79pt}|>{\raggedright}p{247pt}|}
\hline
\multicolumn{2}{|p{326pt}|}{Process dictionary BIFs}\tabularnewline
\hline
\texttt{put(Key, Value)} & Saves the \texttt{Value} under the \texttt{Key} or replaces an older value\tabularnewline
\hline
\texttt{get(Key)} & Retrieves the value stored under \texttt{Key} or \texttt{undefined}\tabularnewline
\hline
\texttt{get()} & Returns the entire process dictionary as a list of \texttt{\{Key, Value\}} terms\tabularnewline
\hline
\texttt{get\_keys(Value)} & Returns a list of keys that have the value \texttt{Value}\tabularnewline
\hline
\texttt{erase(Key)} & Deletes \texttt{\{Key, Value\}}, if any, and returns \texttt{Key}\tabularnewline
\hline
\texttt{erase()} & Returns the entire process dictionary and deletes it\tabularnewline
\hline
\end{tabular}
\end{center}

Process dictionaries could be used to keep global variables within an application,
but the extensive use of them for this is usually regarded as poor programming style.


\chapter{Обработка ошибок}
\label{errorhandling}

Эта глава описывает обработку ошибок внутри процесса. Такие ошибки известны ещё, 
как \textbf{исключения}.


\section{Классы исключений и причины ошибок}

\begin{center}
\begin{tabular}{|>{\raggedright}p{55pt}|>{\raggedright}p{350pt}|}
\hline
\multicolumn{2}{|p{326pt}|}{Классы исключений}\tabularnewline
\hline
\texttt{error} &
Ошибка времени выполнения, например, при применении оператора к недопустимым типам
аргументов. Ошибки времени выполнения могут быть вызваны программистом с помощью
встроенной функции \texttt{erlang:error(Причина)} или 
\texttt{erlang:error(Причина, Аргументы)} \tabularnewline
\hline
\texttt{exit}  &
Процесс вызвал \texttt{exit(Причина)}, см. раздел \ref{processes:termination}
о завершении процессов \tabularnewline
\hline
\texttt{throw}  & 
Процесс вызвал \texttt{throw(Выражение)}, см. раздел 
\ref{errorhandling:catchthrow} о бросках исключений\tabularnewline
\hline
\end{tabular}
\end{center}

Появление исключения приводит к аварийной остановке процесса, то есть его 
исполнение останавливается и он и его данные удаляются из системы. Также это 
действие называется \emph{уничтожением} (termination). После этого сигналы выхода
посылаются всем связанным процессам. Исключение состоит из класса, причины выхода
и копии стека вызовов. Стек вызовов можно сформировать и получить в удобном виде с 
помощью функции \texttt{erlang:get\_stacktrace/0}.

Ошибки времени исполнения и другие исключения могут не приводить к смерти 
процесса, если использовать выражения \texttt{try} и \texttt{catch}.

Для исключений, имеющих класс \texttt{error}, например для обычных ошибок времени 
выполнения, \textbf{причиной выхода} будет кортеж \texttt{\{Причина, Стек\}} где
\texttt{Причина} --- это терм, указывающий более точно на тип ошибки.

\begin{center}
\begin{tabular}{|>{\raggedright}p{110pt}|>{\raggedright}p{320pt}|}
\hline
\multicolumn{2}{|p{326pt}|}{Причины выхода}\tabularnewline
\hline
\texttt{badarg}  & 
Передан аргумент недопустимого типа. \tabularnewline
\hline
\texttt{badarith}  &
Аргумент в арифметическом выражении имеет недопустимый тип. \tabularnewline
\hline
\texttt{\{badmatch, Значение\}}  & 
Вычисление сопоставления с образцом прошло неудачно. \texttt{Значение} не совпало
с образцом.
\tabularnewline
\hline
\texttt{function\_clause}  & 
Не найдено ни одного подходящего уравнения функции по образцам аргументов и
охранных выражений при вычислении вызова функции. \tabularnewline
\hline
\texttt{\{case\_clause, Значение\}}  & 
При вычислении выражения \texttt{case} \texttt{Значение} не совпало ни с чем. \tabularnewline
\hline
\texttt{if\_clause}  & 
Ни одна из веток выражения \texttt{if} не оказалась равна \texttt{true}.
\tabularnewline
\hline
\texttt{\{try\_clause, Value\}}  &
При вычислении секции выражения \texttt{try} со \texttt{Значением} не совпала ни
одна из веток. \tabularnewline
\hline
\texttt{undef}  & 
Функция не была найдена при попытке исполнить вызов функции \tabularnewline
\hline
\texttt{\{badfun, Fun\}}  & 
Что-то не так с анонимной функцией \texttt{Fun} \tabularnewline
\hline
\texttt{\{badarity, Fun\}}  & 
Функция вызвана с неверным количеством аргументов. Значение \texttt{Fun} описывает
и её и переданные аргументы \tabularnewline
\hline
\texttt{timeout\_value}  & 
Значение таймаута в \texttt{receive}$...$\texttt{after} было вычислено и оказалось
не целым 32-битным числом и не атомом \texttt{infinity} \tabularnewline
\hline
\texttt{noproc}  & 
Попытка создать связь с несуществующим процессом \tabularnewline
\hline
\texttt{\{nocatch, Значение\}}  & 
Попытка выполнить \texttt{throw} за пределами кода, защищённого оператором
\texttt{catch}. \texttt{Значение} --- терм, который был брошен \tabularnewline
\hline
\texttt{system\_limit}  & Сработало одно из системных ограничений, заданных 
реализацией виртуальной машины или операционной системы \tabularnewline
\hline
\end{tabular}
\end{center}

\texttt{Стек} --- это цепочка вызовов функций, которые исполнялись в тот момент,
когда произошла ошибка, даётся в виде списка кортежей \texttt{\{Модуль, Имя, 
Арность\}}, где первым идёт самый недавний вызов. Иногда самый последний вызов 
вместо \texttt{Арности} будет содержать \texttt{Аргументы}: \texttt{\{Модуль, 
Имя, Аргументы\}}.


\section{Catch и throw}
\label{errorhandling:catchthrow}
\begin{erlangru}
catch Выражение
\end{erlangru}

Такая запись возвращает результат вычисления \texttt{Выражения} если во время 
вычисления не возникло исключение. В случае исключения возвращаемое значение будет
кортежем с информацией о нём.

\begin{erlangru}
{'EXIT', {Причина, Стек}}
\end{erlangru}

Такое исключение считается \emph{пойманным}. Непойманное исключение приведёт к 
уничтожению процесса. Если исключение было вызвано вызовом функции 
\texttt{exit(Терм)} то будет возвращён кортеж с причиной в виде 
\texttt{\{'EXIT',Терм\}}. Если исключение возникло при вызове \texttt{throw(Терм)},
тогда будет возвращено значение \texttt{Терма}.

\texttt{catch 1+2} \resultingin \texttt{3}\\
\texttt{catch 1+a } \resultingin \texttt{\{'EXIT',\{badarith,[...]\}\}}

Оператор \texttt{catch} имеет низкий приоритет и выражения, использующие его часто
требуют заключения в блок \texttt{begin..end} или круглые скобки.

\texttt{A = (catch 1+2)} \resultingin \texttt{3}

Встроенная функция \texttt{throw(Выражение)} используется для \emph{нелокального}
выхода из функций. Её можно вызывать только под защитой \texttt{catch}, что 
возвратит результат вычисления \texttt{Выражения}.

\texttt{catch begin 1,2,3,throw(four),5,6 end} \resultingin \texttt{four}

Если \texttt{throw/1} вычисляется за пределами оператора \texttt{catch}, то 
возникнет ошибка времени выполнения \texttt{nocatch}.

Оператор \texttt{catch} не спасёт процесс от завершения по сигналу выхода из 
другого связанного с ним процесса (если только не был включен режим перехвата 
сигналов выхода, trap\_exit).


\section{Try}
\label{errorhandling:try}

Выражение \texttt{try} способно различить различные классы исключений.  Следующий
пример эмулирует поведение описанного чуть выше \texttt{catch Выражение}:

\begin{erlang}
try Expr
catch
    throw:Term -> Term;
    exit:Reason -> {'EXIT', Reason};
    error:Reason -> {'EXIT',{Reason, erlang:get_stacktrace()}}
end
\end{erlang}

Полное описание \texttt{try} следующее:

\begin{erlangru}
try Выражение [of
    Образец1 [when ОхранныеВыражения1] -> Тело1;
    ...;
    ОбразецN [when ОхранныеВыраженияN] -> ТелоN]
[catch
    [КлассОшибки1:]ОбразецИсключения1 [when ОхранаИсключения1] ->
        ТелоИсключения1;
    ...;
    [КлассОшибкиN:]ОбразецИсключенияN [when ОхранаИсключенияN] ->
        ТелоИсключенияN]
[after ТелоПосле]
end
\end{erlangru}

В наличии должны обязательно быть как минимум одна строка \texttt{catch} или одно
предложение \texttt{after}.  Может дополнительно использоваться оператор
\texttt{of} после \texttt{Выражения}, который добавляет вычисление \texttt{case}
к значению \texttt{Выражения}.

\texttt{try} возвратит значение \texttt{Выражения}, если только не произойдёт 
исключение во время его вычисления.  Тогда исключение \emph{ловится} и ряд 
\texttt{ОбразцовИсключения} с подходящим \texttt{КлассомОшибки} сопоставляются
один за другим с пойманным исключением.  Если \texttt{КлассОшибки} не указан, то
подразумевается \texttt{throw}. Если сопоставление удачно проходит и необязательные
выражения \texttt{ОхраныИсключения} тоже равны \texttt{true}, то вычисляется 
соответствующее \texttt{ТелоИсключения} и его результат становится возвращаемым
значением.

Если не найдено подходящего \texttt{ОбразцаИсключения} с таким же 
\texttt{КлассомОшибки} и выражениями \texttt{ОхраныИсключения}, равными 
\texttt{true}, то исключение передаётся дальше, как если бы начальное 
\texttt{Выражение} не было заключено в \texttt{try}.  Исключение, происходящее
во время вычисления \texttt{ТелаИсключения} не будет поймано.

Если ни один из \texttt{Образцов} не совпал, то произойдёт ошибка времени 
исполнения \texttt{try\_clause}.

Если определено, то \texttt{ТелоПосле} всегда вычисляется \textbf{после} всех 
операций в\linebreak
\texttt{try..catch}, независимо от возникновения ошибки. Его 
возвращаемое значение игнорируется и не влияет на значение всего оператора 
\texttt{try} (как если бы \texttt{after} не было). \texttt{ТелоПосле} вычисляется
даже если исключение произошло в \texttt{Теле} или \texttt{ТелеИсключения}, в этом
случае исключение передаётся выше по коду.

%% The \texttt{AfterBody} is evaluated after either \texttt{Body} or
%% \texttt{ExceptionBody} no matter which one. The evaluated value of the
%% \texttt{AfterBody} is lost; the return value of the try expression is
%% the same with an after section as without. Even if an exception occurs
%% during evaluation of \texttt{Body} or \texttt{ExceptionBody}, the
%% \texttt{AfterBody} is evaluated. In this case the exception is caught
%% and passed on after the \texttt{AfterBody} has been evaluated, so the
%% exception from the try expression is the same with an after section as
%% without.

Исключение, которое происходит во время вычисления \texttt{ТелаПосле} не ловится,
то есть если \texttt{ТелоПосле} вычислилось по причине исключения в 
\texttt{Выражении}, \texttt{Теле} или \texttt{ТелеИсключения}, то оригинальное
исключение будет потеряно и вместо него полёт вверх по стеку вызовов продолжит уже
новое исключение.


\chapter{Distributed Erlang}
\label{distribution}

A \textbf{distributed Erlang system} consists of a number of Erlang
runtime systems communicating with each other. Each such runtime
system is called a \textbf{node}.  Nodes can reside on the same
host or on different hosts connected through a network.  The
standard distribution mechanism is implemented using TCP/IP sockets
but other mechanisms can also be implemented.

Message passing between processes on different nodes, as well as links
and monitors, is transparent when using \texttt{Pid}s.  However, registered
names are local to each node.  A registered process at a particular
node is referred to as \texttt{\{Name,Node\}}.

The Erlang Port Mapper Daemon \textbf{epmd} is automatically started
on every host where an Erlang node is started.  It is responsible for
mapping the symbolic node names to machine addresses.


\section{Nodes}
A \textbf{node} is an executing Erlang runtime system which has been
given a name, using the command line flag \texttt{-name} (long name)
or \texttt{-sname} (short name).

The format of the node name is an atom \texttt{Name@Host} where
\texttt{Name} is the name given by the user and \texttt{Host} is the
full host name if long names are used, or the first part of the host
name if short names are used. \texttt{node()} returns the name of the
node. Nodes using long names cannot communicate with nodes using
short names.


\section{Node connections}
The nodes in a distributed Erlang system are fully connected. The
first time the name of another node is used, a connection attempt to
that node will be made. If a node A connects to node B, and node B has
a connection to node C, then node A will also try to connect to node
C. This feature can be turned off using the command line flag:

\texttt{~~~~-connect\_all false}

If a node goes down, all connections to that node are removed. The
BIF:

\begin{erlang}
erlang:disconnect(Node)
\end{erlang}

disconnects \texttt{Node}. The BIF \texttt{nodes()} returns the list of
currently connected (visible) nodes.


\section{Hidden nodes}
It is sometimes useful to connect to a node without also connecting to
all other nodes. For this purpose, a \textbf{hidden node} may be
used. A hidden node is a node started with the command line flag
\texttt{-hidden}.  Connections between hidden nodes and other nodes must be set
up explicitly. Hidden nodes do not show up in the list of nodes
returned by \texttt{nodes()}. Instead, \texttt{nodes(hidden)} or \texttt{nodes(connected)} must
be used. A hidden node will not be included in the set of nodes that
the module global keeps track of.

A \textbf{C node} is a C program written to act as a hidden node in a
distributed Erlang system. The library \texttt{erl\_interface}
contains functions for this purpose.


\section{Cookies}
Each node has its own \textbf{magic cookie}, which is an atom. The
Erlang network authentication server (auth) reads the cookie in the
file \texttt{\$HOME/.erlang.cookie}. If the file does not exist, it
will be created with a random string as content.

% FRMB CHECK: the implication here is that "erlang:set_cookie(node(), Cookie)"
% sets *this* node's cookie, as well as the cookie that will be used to connect
% to other nodes, if not explicitly set otherwise.

The permissions of the file must be set to octal 400 (read-only
by user).  The cookie of the local node may also be set using the BIF
\texttt{erlang:set\_cookie(node(), Cookie)}.

The current node is only allowed to communicate with another node
\texttt{Node2} if it knows its cookie.  If this is different from the current node (whose cookie will
be used by default) it must be explicitly set with the BIF \texttt{erlang:set\_cookie(Node2, Cookie2)}.


\section{Distribution BIFs}

\begin{center}
\begin{tabular}{|>{\raggedright}p{156pt}|>{\raggedright}p{170pt}|}
\hline
\multicolumn{2}{|p{326pt}|}{Distribution BIFs}\tabularnewline
\hline
\texttt{node()}  & Returns the name of the current node. Allowed in guards\tabularnewline
\hline
\texttt{is\_alive()}  & Returns true if the runtime system is a node and can connect to
other nodes, false otherwise\tabularnewline
\hline
\texttt{erlang:get\_cookie()}  & Returns the magic cookie of the current node\tabularnewline
\hline
\texttt{set\_cookie(Node, Cookie)} & Sets the magic cookie used when connecting to \texttt{Node}.
If \texttt{Node} is the current node, \texttt{Cookie} will be used when connecting to all new nodes\tabularnewline
\hline
\texttt{nodes()}  & Returns a list of all visible nodes to which the current node is connected
to\tabularnewline
\hline
\texttt{nodes(connected\textbar{}hidden)}  & Returns a list not only of visible nodes,
but also hidden nodes and previously known nodes, etc. \tabularnewline
\hline
\texttt{monitor\_node(Node,}\break\texttt{\phantom{xxxx}true\textbar{}false)}  & Monitors the status of \texttt{Node}. A message
\texttt{\{nodedown, Node\}} is received if the connection to it is lost\tabularnewline
\hline
\texttt{node(Pid\textbar{}Ref\textbar{}Port)}  & Returns the node where the argument is
located\tabularnewline
\hline
\texttt{erlang:disconnect\_node(Node)}  & Forces the disconnection of \texttt{Node}\tabularnewline
\hline
\texttt{spawn[\_link\textbar{}\_opt](Node,}\break\texttt{\phantom{xxxx}Module, Function, Args)}  & Creates a process
at a remote node\tabularnewline
\hline
\texttt{spawn[\_link\textbar{}\_opt](Node, Fun)}  & Creates a process at a remote node\tabularnewline
\hline
\end{tabular}
\end{center}


\section{Distribution command line flags}

\begin{center}
\begin{tabular}{|>{\raggedright}p{102pt}|>{\raggedright}p{224pt}|}
\hline
\multicolumn{2}{|p{326pt}|}{Distribution command line flags}\tabularnewline
\hline
\texttt{-connect\_all false}  & Only explicit connection set-ups will be used\tabularnewline
\hline
\texttt{-hidden}  & Makes a node into a hidden node\tabularnewline
\hline
\texttt{-name Name}  & Makes a runtime system into a node, using long node names\tabularnewline
\hline
\texttt{-setcookie Cookie}  & Same as calling \linebreak{}
\texttt{erlang:set\_cookie(node(), Cookie))}\tabularnewline
\hline
\texttt{-sname Name}  & Makes a runtime system into a node, using short node names\tabularnewline
\hline
\end{tabular}
\end{center}


\section{Distribution modules}
There are several modules available which are useful for distributed programming:

\begin{center}
\begin{tabular}{|>{\raggedright}p{93pt}|>{\raggedright}p{233pt}|}
\hline
\multicolumn{2}{|p{326pt}|}{Kernel modules useful for distribution}\tabularnewline
\hline
\texttt{global}  & A global name registration facility\tabularnewline
\hline
\texttt{global\_group}  & Grouping nodes to global name registration groups\tabularnewline
\hline
\texttt{net\_adm}  & Various net administration routines\tabularnewline
\hline
\texttt{net\_kernel}  & Erlang networking kernel\tabularnewline
\hline
\multicolumn{2}{|p{326pt}|}{STDLIB modules useful for distribution}\tabularnewline
\hline
\texttt{slave}  & Start and control of slave nodes\tabularnewline
\hline
\end{tabular}
\end{center}

\chapter{Порты и драйверы портов}
\label{ports}

\textbf{Порты} предоставляют байто-ориентированный интерфейс к внешним программам и
связывается с процессами Erlang посылая и принимая сообщения в виде списков байтов.
Процесс Erlang, который создаёт порт, называется \textbf{владельцем порта} или
\textbf{подключенным к порту процессом}.  Все коммуникации в и из порта должны
пройти через владельца порта.  Если владелец порта завершает работу, порт тоже
закроется (а также и внешняя программа, подключенная к порту, если она была
правильно написана и среагирует на закрытие ввода/вывода).

Внешняя программа является другим процессом операционной системы. По умолчанию, он
должен считывать данные из стандартного входа (файловый дескриптор 0) и отвечать
на стандартный вывод (файловый дескриптор 1).  Внешняя программа должна завершать
свою работу когда порт закрывается (ввод/вывод закрывается).


\section{Драйверы портов}

Драйверы портов обычно пишутся на языке С и динамически подключаются к системе
исполнения Erlang. Встроенный драйвер ведёт себя как порт и называется 
\textbf{драйвером порта}.  Однако, ошибка в драйвере порта может привести к 
нестабильности во всей системе Erlang, утечкам памяти, зависаниям и краху системы.


\section{Встроенные функции для портов}

\begin{center}
\begin{tabular}{|>{\raggedright}p{140pt}|>{\raggedright}p{300pt}|}
\hline
\multicolumn{2}{|p{440pt}|}{Функция для создания порта} \tabularnewline
\hline
\texttt{open\_port(ИмяПорта, НастройкиПорта)} & 
Возвращает \textbf{идентификатор порта} \texttt{Порт}, как результат открытия 
нового Erlang-порта.  Сообщения могут быть отправлены в и получены через 
идентификатор порта, так же как это можно делать с идентификаторами процессов.
Идентификаторы портов могут быть связаны с процессами, или зарегистрированы под
каким-либо именем с помощью \texttt{link/1} и \texttt{register/2}. \tabularnewline
\hline
\end{tabular}
\end{center}

\texttt{ИмяПорта} обычно является кортежем вида \texttt{\{spawn,Команда\}}, где
строка \texttt{Команда} является именем внешней программы.  Внешняя программа 
выполняется за пределами Erlang-системы, если только не найден драйвер порта с
именем \texttt{Команда}.  Если драйвер найден, он будет активирован вместо команды.

\texttt{НастройкиПорта} --- это список настроек (опций) для порта. Список обычно
содержит как минимум один кортеж \texttt{\{packet,N\}}, указывающий, что данные,
пересылаемые между портом и внешней программой, предваряются N-байтовым индикатором
длины.  Разрешённые значения для \texttt{N} --- 1, 2 или 4.  Если двоичные данные
должны использоваться вместо списков байтов, то должна быть включена опция
\texttt{binary}.

Владелец порта Pid связывается с \texttt{Портом} с помощью отправки и получения
Erlang-сообщений.  (Любой процесс может послать сообщение в порт, но ответы от 
порта всегда будут отправлены только владельцу порта).

\begin{center}
\begin{tabular}{|>{\raggedright}p{140pt}|>{\raggedright}p{300pt}|}
\hline
\multicolumn{2}{|p{440pt}|}{Сообщения, отсылаемые в порт}\tabularnewline
\hline
\texttt{\{Pid, \{command, Данные\}\}}  & 
Посылает \texttt{Данные} в порт. \tabularnewline
\hline
\texttt{\{Pid, close\}}  & 
Закрывает порт. Если порт был открыт, он отвечает сообщением 
\texttt{\{Порт, closed\}}, когда все буферы были сброшены и порт закрылся.
\tabularnewline
\hline
\texttt{\{Pid, \{connect, НовыйPid\}\}}  & 
Устанавливает владельца \texttt{Порта} равным \texttt{НовомуPid}. Если порт был
открыт, он отвечает сообщением \texttt{\{Порт, connected\}} старому владельцу. 
Заметьте, что старый владелец порта остаётся связанным с портом, тогда как новый
--- нет. \tabularnewline
\hline
\end{tabular}
\end{center}

Данные должны быть списком ввода-вывода. Список ввода-вывода --- это либо двоичные
данные, либо смешанный (возможно вложенный) список двоичных данных и целых чисел
в диапазоне от 0 до 255.

\begin{center}
\begin{tabular}{|>{\raggedright}p{140pt}|>{\raggedright}p{300pt}|}
\hline
\multicolumn{2}{|p{440pt}|}{Сообщения, получаемые из порта}\tabularnewline
\hline
\texttt{\{Порт, \{data, Данные\}\}}  & 
Данные получены от внешней программы \tabularnewline
\hline
\texttt{\{Порт, closed\}}  & 
Ответ на команду \texttt{Порт ! \{Pid,close\}} \tabularnewline
\hline
\texttt{\{Порт, connected\}}  & 
Ответ на команду \texttt{Порт ! \{Pid,\{connect, NewPid\}\}} \tabularnewline
\hline
\texttt{\{'EXIT', Порт, Причина\}}  &
Присылается, если порт был отключен по какой-либо причине. \tabularnewline
\hline
\end{tabular}
\end{center}

Вместо того, чтобы отправлять и получать сообщения, имеется ряд встроенных функций,
которые можно использовать.  Они могут быть вызваны любым процессом, а не только
владельцем порта.

\begin{center}
\begin{tabular}{|>{\raggedright}p{140pt}|>{\raggedright}p{300pt}|}
\hline
\multicolumn{2}{|p{440pt}|}{Встроенные функции для работы с портами}\tabularnewline
\hline
\texttt{port\_command(Порт, Данные)}  & 
Отправляет \texttt{Данные} в \texttt{Порт} \tabularnewline
\hline
\texttt{port\_close(Порт)}  & 
Закрывает \texttt{Порт} \tabularnewline
\hline
\texttt{port\_connect(Порт, НовыйPid)}  & 
Устанавливает владельца \texttt{Порта} равным \texttt{НовомуPid}. Старый владелец
остаётся связанным с портом и должен сам вызвать \texttt{unlink(Порт)} если связь
не требуется. \tabularnewline
\hline
\texttt{erlang:port\_info(\\
Порт, Элемент)}  & 
Возвращает информацию о \texttt{Порте} с ключом \texttt{Элемент} \tabularnewline
\hline
\texttt{erlang:ports()}  & 
Возвращает список всех открытых портов на текущем узле \tabularnewline
\hline
\end{tabular}
\end{center}

Есть несколько дополнительных встроенных функций, которые применимы только к 
драйверам портов: это \texttt{port\_control/3} и \texttt{erlang:port\_call/3}.
\chapter{Загрузка кода}
\label{code}

Erlang поддерживает обновление кода во время работы без остановки системы.  
Замена кода выполняется на уровне модулей.

Код модуля может существовать в системе в двух версиях: \textbf{текущая} и
\textbf{старая}. Когда модуль загружается в систему в первый раз, код становится
\emph{текущим}.  Если загружается новая версия модуля, то код предыдущей версии,
уже имеющийся в памяти, становится \emph{старым} и новая загруженная версия 
становится \emph{текущей}.  Обычно модуль автоматически загружается, когда 
вызвана одна из находящихся в нём функций.  Если модуль уже загружен, то он
должен быть явно перезагружен в новую версию.

И старый и текущий код полностью функциональны и могут использоваться 
одновременно. Полностью определённые вызовы функций (с именем модуля) всегда 
ссылаются на текущую версию кода.  Однако старый код может продолжать 
исполняться другими процессами.

Если загрузить третью версию загруженного модуля, то сервер кода удалит 
(операция называется \emph{purge}) старый код и все процессы, всё ещё 
использующие его, будут принудительно завершены.  Затем третья версия 
становится \emph{текущей} и предыдущий текущий код становится \emph{старым}.

Чтобы перейти от старого кода к текущему, процесс должен выполнить один 
полностью определённый вызов функции (с именем модуля).

\pagebreak
\begin{erlang}
-module(mod).
-export([loop/0]).

loop() ->
    receive
        code_switch ->
            mod:loop();
        Msg ->
            ...
            loop()
    end.
\end{erlang}

Чтобы заставить процесс (в этом примере) сменить версию кода, отправьте ему
сообщение \texttt{code\_switch}.  Процесс после этого выполнит полностью 
определённый вызов \texttt{mod:loop()} и это переключит его на текущую версию 
кода. Заметьте, что \texttt{mod:loop/0} должна быть для этого экспортирована.
\chapter{Macros}
\label{macros}

\section{Defining and using macros}

\begin{erlang}
-define(Const, Replacement).
-define(Func(Var1, ..., VarN), Replacement).
\end{erlang}

A \textbf{macro} must be defined before it is used but a macro
definition may be placed anywhere among the attributes and function
declarations of a module. If a macro is used in several modules it is
advisable to put the macro definition in an include file. A macro is
used as follows:

\begin{erlang}
?Const
?Func(Arg1,...,ArgN)
\end{erlang}

Macros are expanded during compilation. A macro reference
\texttt{?Const} is replaced by \texttt{Replacement} like this:

\begin{erlang}
-define(TIMEOUT, 200).
...
call(Request) ->
    server:call(refserver, Request, ?TIMEOUT).
\end{erlang}

is expanded to:

\begin{erlang}
call(Request) ->
    server:call(refserver, Request, 200).
\end{erlang}

A macro reference \texttt{?Func(Arg1, ..., ArgN)} will be replaced by
\texttt{Replacement}, where all occurrences of a variable \texttt{VarX}
from the macro definition are replaced by the corresponding argument
\texttt{ArgX}.

\begin{erlang}
-define(MACRO1(X, Y), {a, X, b, Y}).
...
bar(X) ->
    ?MACRO1(a, b),
    ?MACRO1(X, 123).
\end{erlang}

will be expanded to:

\begin{erlang}
bar(X) ->
    {a, a, b, b},
    {a, X, b, 123}.
\end{erlang}

To view the result of macro expansion, a module can be compiled with
the \texttt{`P'} option:

\begin{erlang}
compile:file(File, ['P']).
\end{erlang}

This produces a listing of the parsed code after preprocessing and
parse transforms, in the file \texttt{File.P}.


\section{Predefined macros}

\begin{center}
\begin{tabular}{|>{\raggedright}p{103pt}|>{\raggedright}p{223pt}|}
\hline
\multicolumn{2}{|p{326pt}|}{P{\large{}redefined} macros}\tabularnewline
\hline
\texttt{?MODULE} & The name of the current module\tabularnewline
\hline
\texttt{?MODULE\_STRING} & The name of the current module, as a string\tabularnewline
\hline
\texttt{?FILE} & The file name of the current module\tabularnewline
\hline
\texttt{?LINE} & The current line number\tabularnewline
\hline
\texttt{?MACHINE} & The machine name, 'BEAM'\tabularnewline
\hline
\end{tabular}
\end{center}


\section{Flow Control in Macros}

\begin{erlang}
-undef(Macro).      % This inhibits the macro definition.

-ifdef(Macro).
    %% Lines that are evaluated if Macro was defined
-else.
    %% If the condition was false, these lines are evaluated instead.
-endif.
\end{erlang}

\texttt{ifndef(Macro)} can be used instead of \texttt{ifdef} and means
the opposite.

\begin{erlang}
-ifdef(debug).
-define(LOG(X), io:format("{~p,~p}:~p~n",[?MODULE,?LINE,X])).
-else.
-define(LOG(X), true).
-endif.
\end{erlang}

If \texttt{debug} is defined when the module is compiled,
\texttt{?LOG(Arg)} will expand to a call to \texttt{io:format/2} and
provide the user with some simple trace output.


\section{Stringifying Macro Arguments}
\texttt{??Arg}, where \texttt{Arg} is a macro argument expands to the
argument in the form of a string.

\begin{erlang}
-define(TESTCALL(Call), io:format("Call ~s: ~w~n", [??Call, Call])).

?TESTCALL(myfunction(1,2)),
?TESTCALL(you:function(2,1)),
\end{erlang}

results in:

\begin{erlang}
io:format("Call ~s: ~w~n", ["myfunction(1,2)", m:myfunction(1,2)]),
io:format("Call ~s: ~w~n", ["you:function(2,1)", you:function(2,1)]),
\end{erlang}

That is, a trace output with both the function called and the
resulting value.

\chapter{Further Reading and Resources}

Following websites provide in-depth explanation of topics and concepts briefly covered in this document:

\begin{itemize}
	\item Official Erlang documentation: \url{http://www.erlang.org/doc/}
	\item Learn You Some Erlang for Great Good: \url{http://learnyousomeerlang.com/}
	\item Tutorials section at Erlang Central: \url{https://erlangcentral.org/wiki/index.php?title=Category:HowTo}
\end{itemize}

Still have questions? erlang-questions mailing list (\url{http://erlang.org/mailman/listinfo/erlang-questions}) is a good place for general discussions about Erlang/OTP, the language, implementation, usage and beginners questions.



\end{document}
